\documentclass[UTF8,11pt]{ctexart}
% Geometry mirrors monograph.tex exactly (see CONVENTIONS.md section 5): the 1in margin box
% is preserved so artifact_qc.py measures the same edges in both editions.
\usepackage[letterpaper,left=72pt,right=72pt,textheight=648pt]{geometry}
\usepackage{amsmath,amssymb,amsthm,mathtools}
\usepackage{booktabs}
\usepackage{enumitem,graphicx,array,longtable}
\usepackage{xcolor}
\usepackage[T1]{fontenc}
\usepackage{hyperref}
\hypersetup{colorlinks=true,linkcolor=blue,urlcolor=blue,citecolor=blue}
% 断行设置。原来的 \tolerance=3000 + \emergencystretch=3em 留下 45 行松行，其中 33 行
% 为最差 badness 10000；原因是 \tolerance 限制了 TeX 的尝试强度，\emergencystretch 无从生效。
% 提到 9999 后 TeX 宁可取略松的行，余量由下面的伸缩量吸收；\hbadness 则不再报告只是略松的行。
\tolerance=9999
\emergencystretch=3em
\hbadness=2500
% CJK has no hyphenation, so a looser line is preferable to a break inside a long \texttt path.
\sloppy
\newtheorem{theorem}{定理}
\newtheorem{proposition}{命题}
\newtheorem{lemma}{引理}
\newtheorem{corollary}{推论}
\theoremstyle{definition}
\newtheorem{definition}{定义}
\newtheorem{construction}{构造}
\theoremstyle{remark}
\newtheorem{remark}{说明}
\newcommand{\D}{\mathbb{D}}
\newcommand{\R}{\mathbb{R}}
\newcommand{\Z}{\mathbb{Z}}
\newcommand{\Q}{\mathbb{Q}}
\newcommand{\Bd}{\partial}
\newcommand{\area}{\operatorname{area}}
\newcommand{\conv}{\operatorname{conv}}
\newcommand{\eps}{\varepsilon}
\newcommand{\rD}{r_{D}}
\newcommand{\dd}{\mathrm{d}}
\newcommand{\half}{\tfrac12}
% 每个部分重新开始小节编号，因此必须给 hyperref 目标加上部分前缀，否则锚点会冲突。
% 同一计数器也作为小节编号前缀，使「第 IV 部分 4.1 节」这类引用无歧义。
\newcounter{partanchor}
\renewcommand{\thesection}{\thepartanchor.\arabic{section}}
\renewcommand{\theHsection}{p\thepartanchor.\arabic{section}}
% 三级编号（"1.4.5"）加两位页码放不进默认的 \@pnumwidth，曾使目录某行溢出 5.96pt。
% 加宽页码框并给引导点留余量；只影响目录，不改变正文版心。
\makeatletter
\renewcommand{\@pnumwidth}{2.2em}
\renewcommand{\@tocrmarg}{3.2em}
\makeatother
% 字体取舍（有意为之，勿当 bug 修）：离线 Tectonic bundle 缺粗体数学斜体 cmmib8 与部分 8pt 字号，
% 故拉丁向量用 \mathbf，希腊向量用普通斜体希腊字母。详见 CONVENTIONS.md 第 5 节。
\title{\bf 用100个等圆盘覆盖单位圆盘\\[2mm]
\large 全部路线的系统讲解，以及当前的上下界}
\author{公开计算研究档案\\\small \texttt{springoxn/disk-covering-100}}
\date{2026年10月}
\begin{document}
\maketitle

\begin{abstract}
\noindent
本文研究 $r_{D}(100)$：用 $100$ 个全等闭圆盘覆盖单位闭圆盘 $D$ 所需的最小半径。本专著按路线逐条、
系统地讲解整个研究纲领：试过的每一种构造、闭合的每一种松弛、以及被证伪或仍未决的每一条路线。
当前可机器重放的严格结论是
\[
\boxed{\;0.113235\;<\;\rD(100)\;\le\;0.1154469548855829\;}
\]
下界是有理半径 $R=22647/200000$ 处的六分支排除证书：$B31$--$B35$ 由半径单调性从
$R=2831/25000=0.11324$ 的证书继承，而尾分支 $B\ge36$ 在 $R$ 处直接对
$1\,841\,089$ 个有理箱重放。上界是一个显式重放的 $100$ 中心构型。两者都不给出精确最优值。
邻近目标 $0.11324$ 仍未闭合：在冻结的 capped-pair 常数下，分支 $B=36$ 存在一个精确有理障碍点，
该点处局部 pair 松弛超过 $B=36$ 面积预算所能支付的截距。本文说明为什么该障碍无法靠增加箱预算或
提高截距消除，并记录未决区间 $(0.113235,\,0.11324]$。

下文每一个定量断言都绑定到 \texttt{verify/out/} 中的机器可读产物。本文刻意保留失败：
一份只含成功的研究记录，无法告诉下一个读者哪些松弛太弱。
\end{abstract}

\tableofcontents

\section*{前言：本文档是什么，以及如何阅读}
\addcontentsline{toc}{section}{前言}

本专著不是一条定理加一个证明。它是一项研究纲领的地图。
用 $100$ 个等圆盘覆盖单位圆盘是一个连续、非凸的问题，其组合类型
随半径而变，没有已知方法能一步求解。实际存在的是一组 \emph{路线}：
产生上界的构造、产生下界的松弛，以及大量被尝试过的变体——它们要么
闭合了某个子情形，要么产生了一份诊断，要么被证伪。本文档的目的是
系统地呈现这些路线，使读者能够看到 (i) 当前的前沿是什么，
(ii) 哪个不等式承载了其中的哪一部分，以及 (iii) 剩余间隙究竟位于何处。

\subsection*{一页纸的现状综述}

\begin{center}
\footnotesize
\begin{tabular}{@{}>{\raggedright\arraybackslash}p{0.30\textwidth}>{\raggedright\arraybackslash}p{0.60\textwidth}@{}}
\toprule
条目 & 数值 / 状态 \\
\midrule
已接受的严格区间 &
$0.113235<\rD(100)\le 0.1154469548855829$ \\[1mm]
已接受的下界（精确有理数） &
$R_\star=22647/200000=0.113235$ \\[1mm]
已接受的上界（精确有理数） &
$U_\star=1154469548855829/10^{16}$ \\[1mm]
相对间隙 $(U_\star-R_\star)/U_\star$ &
约 $1.92\%$（仅作展示） \\[1mm]
下界引擎 &
六个互斥的边界计数分支（branch）
$B31,B32,B33,B34,B35,B{\ge}36$；$B31$--$B35$ 由半径单调性（monotonicity inheritance）从 $R=2831/25000$ 继承，$B{\ge}36$ 在 $R_\star$ 处重放 \\[1mm]
最弱分支 &
$B\ge 36$，具体是单个取值 $B=36$ \\[1mm]
未达到的目标 &
$0.11324$；冻结的 capped-pair 松弛（capped-pair relaxation）在
$x=y=9921033/12500000$ 处有一个精确障碍点（obstruction point） \\[1mm]
未决窗口 &
$(0.113235,\,0.11324]$；有理候选 $0.113236$ 与 $0.113237$ 已启动并停止，
未产生证书（certificate） \\[1mm]
精确最优值、极小多项式 &
未获得；\texttt{global\_\allowbreak{}optimality\_\allowbreak{}proved = false} \\[1mm]
最佳普适（不依赖构型）下界 &
$2787/25000=0.11148$（长弦束），而更早的
$554570693/5000000000=0.1109141386$ 与
$1111720153/10^{10}=0.1111720153$ 仍然有效且被支配 \\
\bottomrule
\end{tabular}
\end{center}

\subsection*{各部分如何衔接}

\begin{itemize}[leftmargin=1.4em,itemsep=2pt]
\item \textbf{第 I 部分} 固定问题、约定（闭圆盘、中心不受限、有理判定层）与证据纪律。请先读这一部分：
\emph{证书}、\emph{诊断}（diagnostic）、\emph{解析}（analytic）与\emph{继承}（inherited）
这些词在这里有精确含义。
\item \textbf{第 II 部分} 汇集上界路线：有限检查如何替代连续的覆盖命题、带 $\alpha=1/\sqrt{67}$ 的
三角格点 $97{+}3$ 构型、被重放的外部 $100$ 中心构型、协作搜索路线，
以及局部代数（KKT）路线。
\item \textbf{第 III 部分} 汇集不依赖构型的解析下界：面积界
$1/10$、带四条几何引理的整数边计数路线、锐化包络（sharpened envelope）、
长弦二次亏损（long-chord quadratic deficit），以及对周界缺陷（$\lambda$）律的否证。
\item \textbf{第 IV 部分} 是技术核心：六分支引擎。它解释了同半径规范化（normalized）、
计数桥，随后是每个分支 $B31$--$B35$ 与尾分支 $B\ge36$，
包括 capped-pair 松弛及其有理方形区域细分。
\item \textbf{第 V 部分} 是前沿：已接受的界如何从 $0.11323$ 移到 $0.1132301$
再到 $0.113235$，$0.11324$ 为何被阻塞，以及一条真正新的不等式必须做什么。
\item \textbf{第 VI 部分} 是路线目录：每个路线族及其构造、状态与证据路径，
包括被证伪与零产出的那些。
\item \textbf{第 VII 部分} 给出复现命令、外部依赖，以及从第三方材料中吸收内容的边界。
\item \textbf{第 VIII 部分} 列出未决义务，不加任何软化。
\item \textbf{附录} 是证据账本：断言 $\leftrightarrow$ 产物，外加 JSON 字段词表
与 $B=36$ 障碍的精确算术。
\end{itemize}

\subsection*{全文使用的证据等级}

\begin{center}
\footnotesize
\begin{tabular}{@{}c>{\raggedright\arraybackslash}p{0.50\textwidth}>{\raggedright\arraybackslash}p{0.20\textwidth}@{}}
\toprule
等级 & 含义 & 可信度 \\
\midrule
A & 结构化机器产物：带状态、计数、哈希、退出码的 JSON 证书 & 最强 \\
B & 独立重放：由不同脚本或阶段复查同一子断言 & 强 \\
C & 发现层（discovery layer）材料：浮点搜索、候选排序、失败的实验 & 仅用于指导下一步 \\
D & 继承推理，例如把较大半径处的排除用于较小半径 & 仅当组合器显式记录时有效 \\
\bottomrule
\end{tabular}
\end{center}

另有三个限定词不断出现，绝不能与等级混淆：

\begin{itemize}[leftmargin=1.4em,itemsep=2pt]
\item \emph{解析} --- 以文字论证证明、可由人核查，但未在证明助手中形式化。本文档中每一条
「普适」陈述都是解析论证加上精确有限算术；没有一条是内核验证的定理。
\item \emph{\texttt{machine\_\allowbreak{}verified\_\allowbreak{}geometry = false}} --- 有限的有理/区间不等式已被机器检查；从几何到这些不等式的归约没有。
\item \emph{诊断} --- 该计算刻意报告一次失败、一个瓶颈或一个零产出结果，
而\emph{不是}某个分支已闭合的证书。
\end{itemize}

\subsection*{记号}

\begin{center}
\footnotesize
\begin{tabular}{@{}ll@{}}
\toprule
符号 & 含义 \\
\midrule
$D$ & 闭单位圆盘 $\{x\in\R^2:\|x\|\le 1\}$ \\
$B(c,r)$ & 以 $c$ 为圆心、半径 $r$ 的闭圆盘 \\
$\rD(n)$ & $\inf\{r: D\subseteq\bigcup_{i=1}^{n}B(c_i,r)\}$ \\
$R$ & 用于矛盾论证的固定有理 \emph{候选}半径 \\
$U$ & 已认证的上界（一个显式覆盖构型的半径） \\
$B$ & 覆盖所服务的边界弧（boundary arc）（等价地边界顶点）数目 \\
$M$ & 规范化的公共三角剖分的内部面数目 \\
$N$, $m$ & 互异中心的数目，以及有效使用圆盘的数目 \\
$k$ & 互异的服务边界圆盘的数目 \\
$e_i$ & $m_i+a_i$，即被裁剪胞 $Q_i$ 的名义边数 \\
$\Gamma$ & 角度记账中使用的标记内角和总计 \\
$h$ & 严格深顶点的总度数超出量 \\
$C$ & 核心（core）边界的弦数 \\
$\lambda,\mu,L,K$ & 尾松弛的罚系数（penalty coefficient）、斜率（slope）、截距（intercept）与标记常数 \\
$\Phi$ & capped-pair 松弛的局部 pair 损失 \\
\bottomrule
\end{tabular}
\end{center}

\subsection*{关于如何阅读这些数字的警告}

本文档中的每个有理数都按其在证书或报告中出现的样子引用。凡为定位而给出的
十进制数，都标注为展示。凡两个阶段给出略有差异的辅助常数（例如不同半径处
规范化的中心范数平方界），二者都连同各自的来源一并引用，而不做静默调和。
凡一次运行被启动后停止，就报告为已停止 --- 既不报告为失败，也不报告为证明。

\part*{第 I 部分 问题、约定与证据纪律}
\addcontentsline{toc}{part}{第 I 部分 问题、约定与证据纪律}
\stepcounter{partanchor}\setcounter{section}{0}

\section{问题的陈述}

\begin{definition}[覆盖半径（covering radius）与极值量]
令 $D=\{x\in\R^2:\|x\|_2\le1\}$。对中心元组 $C=(c_1,\dots,c_{100})\in(\R^2)^{100}$，置
\[
R_D(C)=\max_{x\in D}\ \min_{1\le i\le 100}\|x-c_i\|_2,
\qquad
\rD(100)=\inf_{C\in(\R^2)^{100}}R_D(C).
\]
覆盖圆盘是\emph{闭}的；它们可以重叠，也可以伸出 $D$ 之外；中心不受限于 $D$。
本纲领的任务是确定 $\rD(100)$，并给出它的一个显式极小多项式、一个达到它的
显式构型、一个覆盖性证明与一个全局最优性证明，全部由可复现的计算机辅助材料支撑。
\end{definition}

五项形式化要求如下，按其约束工作的顺序排列：
\begin{enumerate}[leftmargin=1.6em,itemsep=2pt]
\item \textbf{极小多项式与隔离。} 列出最优半径的本原整系数多项式 $P$ 的每一个系数，
证明其在 $\Q$ 上不可约，并给出有理数 $a<b$，使 $P(a)P(b)\ne0$
且 $P$ 在 $(a,b)$ 内恰有一个根。
\item \textbf{显式构型。} 由精确代数数据给出的一个确定的有序 $100$ 元中心组。
\item \textbf{覆盖性。} 严格证明每个 $x\in D$ 都落在某个中心的 $\alpha$ 邻域内；若把连续命题
归约为有限多次检查，则该归约本身必须被证明是充分的。
\item \textbf{全局最优性。} 对\emph{每个} $C$ 都有 $R_D(C)\ge\alpha$，不得对对称性、接触结构或
组合类型施加未经证明的限制，并且必须显式处理退化情形（中心重合、圆盘冗余、
中心共线、相切接触）。
\item \textbf{计算机辅助证明。} 完整的代码、输入、证书、环境与实际验证结果，
并解释这些结果为何蕴含相应的数学陈述。
\end{enumerate}

\section{问题为何困难}

\subsection{连续、非凸、组合不稳定}

目标函数 $C\mapsto R_D(C)$ 是距离的极小值的极大值：连续、分段光滑且非凸。
其极小化点不由单一光滑系统刻画，因为诱导分解的组合类型随 $r$ 改变：
哪个中心占有一条边界弧（boundary arc）、哪三个中心在某个 Voronoi 顶点相交、
边界上承载多少条弧，全都是构型的函数。因此数值优化器永远只能给出上界。
下界必须排除\emph{所有}构型，而做到这一点的标准途径是转到可能分解的有限
组合目录，再用精确不等式排除目录中的每一项。

\subsection{对称性不是最优性的必要条件}

经典反例是 $n=5$：完全对称的排布给出
$r=\varphi-1=0.6180340\ldots$，而真正的最优值是 $0.6093829\ldots$，
由一个只有单条镜像轴的构型达到。这就是 Connelly 问题（最优覆盖构型是否必须
具有 $n$ 重旋转对称性？），Csik\'os 对 $n=5$ 给出否定回答，对 $n=3$ 给出肯定回答。
对本纲领而言，实际推论很直白：\emph{不}可以把范围限制在对称族内，
然后把结果称为 $\rD(100)$。

\subsection{在 $n=100$ 处枚举不可行}

严格的计算机辅助最优性证明的成熟范式是\emph{组合分类加上区间 Newton /
分支定界}。Nurmela 与 \"Osterg{\aa}rd 把正方形内等圆\emph{装填}的严格最优性
推进到 $n\le27$；他们自己的记述表明，初始组合计数增长极快，
$23\le n\le26$ 花费了以 CPU 月计的时间，而 $n=27$ 需要更换算法。
本纲领对 $n=100$ 的估计与此一致：覆盖构型的轨道数是天文数字。
因此这里的枚举与分类只作为分支内部的\emph{局部}手段使用，绝不作为全局方法。

\subsection{文献中已知的内容（2026 年十月快照）}

单位换算：文献通常用 $\rho(n)$ 表示 $n$ 个单位圆盘能够覆盖一个圆盘时的
最大半径；于是 $\rD(n)=1/\rho(n)$。

\begin{center}
\footnotesize
\begin{tabular}{@{}>{\raggedright\arraybackslash}p{0.20\textwidth}>{\raggedright\arraybackslash}p{0.22\textwidth}>{\raggedright\arraybackslash}p{0.46\textwidth}@{}}
\toprule
$n$ & 状态 & 来源 \\
\midrule
$1$--$4$, $7$ & 精确闭式，最优性平凡 & --- \\
$5$, $6$ & 精确闭式，最优性已证明 & K. Bezdek 1979/1983 \\
$8$, $9$, $10$ & $\rD(n)=1/(1+2\cos\frac{2\pi}{n-1})$，最优性已证明 &
G.~Fejes T\'oth 1996/2005 \\
$11$--$20$ & 已找到构型；某一研究组声称最优，但仍被一个经过整理的数据库
列为未知 & Melissen 1997 / Tan 2018；Friedman 的
\texttt{circovcir} 页面 vs.\ \texttt{leapsinbounds} \\
$21$--$25$ & 已找到构型，未声称最优性证明 & Tan 2018 \\
$26$--$99$ & 没有系统的公开表格 & --- \\
$100$ & 除格点型构造外没有公开数据 & --- \\
\bottomrule
\end{tabular}
\end{center}

那份快照中的三个事实塑造了下面的一切。第一，渐近密度陈述是严格的，但它是渐近的：
$\lim_{\eps\to0^+}\eps^2N(\eps)=2\pi/(3\sqrt3)$，即覆盖密度
$\delta=2\pi/(3\sqrt3)=1.2091996\ldots$（Kershner 1939；Verblunsky 1949）。
要把渐近陈述转成有限 $n$ 的界，需要一个边界修正项，而困难恰恰就在这个修正项上。
第二，唯一成功的「强制对称」范式（Fejes T\'oth 对 $n=8,9,10$ 的证明）只在
「一个中心圆盘加一个环」的结构中有效；$n=25$（中心加三个八边形环）与
$n=100$（许多环）都没有这样的定理。第三，L.~Fejes T\'oth 带周长项的覆盖密度不等式，
及其装填对偶（Oler/Thue--Groemer，$\frac{1}{2\sqrt3}A(C)+\frac14P(C)+1\ge n$），
是「逐胞加边界修正」路线的文献原型 --- 第 III 部分表明它在这里\emph{不够}强，
第 IV 部分则用一个计数加预算的引擎取代它。

\section{约定与规范化}

\begin{enumerate}[leftmargin=1.6em,itemsep=3pt]
\item \textbf{一切都是闭的。} 圆盘、单位圆盘以及所有有理区间都是闭的。
用作外包围的闭区间 $[a,b]$ 意味着真值落在 $[a,b]$ 内；半径相等的闭情形
被保留而不是丢弃。
\item \textbf{判定层。} 每一条被接受的不等式都由整数、\texttt{Fraction}，或由向外舍入
（outward rounding）的有理区间算术（rational interval arithmetic）判定。浮点只出现在发现层：
Newton 根、候选构型、因子筛、LLL/PSLQ 候选关系。发现层的结果必须被精确地
重新判定之后，才能进入证书。
\item \textbf{投影。} 若 $c\notin D$，令 $\Pi(c)$ 为其到凸集 $D$ 上的投影。由于 $\Pi$ 是非扩张的，
$B(\Pi(c),r)\supseteq B(c,r)\cap D$；把每个中心替换为其投影，既保持覆盖性也保持半径。
因此可以假设 $\|c_i\|\le1$。这\emph{不}把中心置于第 III 部分所用的内接多边形 $P$ 之内，
而那里的几何引理之一正是为了去掉这个更早的假设。
\item \textbf{去重。} 若 $i\ne j$ 时 $c_i=c_j$，删去一个圆盘：覆盖性与半径不变，
而中心集缩小。去重之后所有中心互异，且每个 Voronoi 胞具有正面积。
\item \textbf{输出不覆盖。} 证书脚本拒绝覆盖已存在的 JSON 目标。因此重跑一个阶段
需要一个新的输出路径，仓库正是以此区分「这份证书由本次运行产生」与
「存在一个同名的文件」。
\end{enumerate}

\section{证据纪律}

\subsection{绝不能混为一谈的五个层次}

本纲领自己的综合报告区分了五个验证层次，并坚持较早的层次绝不替代较晚的层次：

\begin{center}
\footnotesize
\begin{tabular}{@{}>{\raggedright\arraybackslash}p{0.26\textwidth}>{\raggedright\arraybackslash}p{0.62\textwidth}@{}}
\toprule
层次 & 它所确立的内容 \\
\midrule
1. 文件存在性 / 哈希 & 本次运行实际消费了哪些字节 \\
2. 元数据 / 条目验证 & 拒绝重复节点、错误定义域、未授权
条目，以及无依据的预算添加 \\
3. 精确算术重放 & 对\emph{给定}定义域重新计算必要条件与证书 \\
4. 解析 / 几何桥 & 每个真正的覆盖构型都落入该定义域，
即那些必要约束确实成立 \\
5. 全局覆盖 & 所有分支、图与标号都被枚举或包围 \\
\bottomrule
\end{tabular}
\end{center}

\subsection{\texttt{PASS} 意味着「在 scope 字段范围内通过」}

一条存储的 \texttt{PASS} 恰好断言该证书 \texttt{scope} 字段所说的内容，不多不少。
本纲领中的三个具体事件说明了这一点：

\begin{itemize}[leftmargin=1.4em,itemsep=3pt]
\item \textbf{一个独立脚本复制了同一个错误。} 在阶段 45/46 中，一份边界环带
证书使用了错误的内面过滤器，产生了空的待处理集合；「全局无增益」的接受被判无效。
修正后的运行实际检查了 $32$ 个拥有边界的面、$4776$ 个构型与 $152\,832$ 次精确
环带比较，并新增 $0$ 条排除。状态文件中记录的教训是：那个独立脚本重复了
同一个过滤器错误，因此\emph{退出码 0 不足以证明数学覆盖定义域是正确的}。
\item \textbf{一个位于不同半径的 \texttt{PASS}。} $B31$ 证书是在 $R=2831/25000$ 处建立的；
它只有在借助半径单调性的情况下才能用于更小的半径，而组合器必须显式记录
这一继承。一个半径处的 \texttt{PASS} 不是另一个半径处的 \texttt{PASS}。
\item \textbf{一个 \texttt{PASS\_\allowbreak{}PREMISES}。} $B35$ 诊断证书返回
\texttt{status = PASS\_\allowbreak{}PREMISES} 与 \texttt{B35\_\allowbreak{}closed = false}：所有前提（premises）都被验证，
而该分支\emph{没有}闭合。这个状态词的存在，使得这样一次运行
永远不会被读作闭合。
\end{itemize}

\subsection{负控制是证书的一部分}

几乎每个证书脚本都附带必须被拒绝的破坏性测试：放大的目标圆盘、针对放大圆盘变体的
一个显式未覆盖点、缩小的覆盖半径、放大的覆盖半径、错误的度量、旋转或膨胀的
Voronoi 六边形、被截断的邻居搜索、错误归属的邻居、被移除的相切中心、
被移除的内部中心、把 $97$ 中心构型当作 $100$ 点断言提交、断开的边配对、
破损的面置换、被声明为不活跃的支撑、用 \texttt{python -\allowbreak{}O} 禁用的断言，
以及若干刻意损坏的有理载荷。负控制没有被全部拒绝的运行不被接受。另外两个入口点
还拒绝在 \texttt{python -\allowbreak{}O} 下运行，并报告 \emph{Assertions must remain enabled}。

\subsection{状态词}

\begin{center}
\footnotesize
\begin{tabular}{@{}>{\raggedright\arraybackslash}p{0.26\textwidth}>{\raggedright\arraybackslash}p{0.62\textwidth}@{}}
\toprule
状态 & 含义 \\
\midrule
\texttt{PASS} & 所断言的检查在声明的 scope 内通过 \\
\texttt{PASS\_\allowbreak{}PREMISES} & 前提已验证，\emph{未}声称闭合 \\
\texttt{NOT\_\allowbreak{}CLOSED} & 该分支或目标仍开放 \\
\texttt{EXCLUDED} & 声明的定义域为空（一次真正的排除） \\
\texttt{NOT\_\allowbreak{}EXCLUDED} / \texttt{UNDECIDED} & 非空或未解决；绝不计为已排除 \\
\texttt{capped} & 传播在未达到不动点的情况下触及轮次上限；予以保留，不做判定 \\
\texttt{diagnostic} & 被记录的失败或瓶颈，不是证书 \\
\bottomrule
\end{tabular}
\end{center}

\subsection{绝不能推断出的内容}

证据账本明确陈述了这些禁令，它们正是本文档保留其失败记录可见的原因：
\begin{itemize}[leftmargin=1.4em,itemsep=2pt]
\item 一个通过的局部盒不是全局下界；
\item 一个未被阻塞的障碍点（obstruction point）不是已闭合的分支；
\item 一个存在的构型不是最优性证明；
\item 一次数值上的「通过」打印不是证明；
\item 检查有限多个采样点、检查边界或检查一幅数值图像，都不是对整个闭单位圆盘
已被覆盖的证明。
\end{itemize}

\part*{第二部分：上界——构造及其证书}
\addcontentsline{toc}{part}{第二部分：上界——构造及其证书}
\stepcounter{partanchor}\setcounter{section}{0}

\section{上界路线必须做什么}

上界是一个构造：给出 $100$ 个圆心 $c_1,\dots,c_{100}$ 与一个有理数 $U$，并证明
$\max_{x\in D}\min_i\|x-c_i\|\le U$。因此其数学内容永远是同一个两步归约
\[
\underbrace{D\subseteq\bigcup_i B(c_i,U)}_{\text{continuous statement}}
\;\Longleftarrow\;
\underbrace{\text{finite list of exact box/leaf checks}}_{\text{machine decision}} ,
\]
而全部困难都在这条蕴含式里。本纲领用到三种不同的有限归约，它们值得分开讨论，
因为三者的长处各不相同。

\subsection{归约 A：四叉树／箱细分，配精确距离判据}

用轴平行箱覆盖 $D$。对箱 $X$ 与圆心 $c$，$\|x-c\|^2$ 在 $X$ 上的区间赋值给出
$[\ell(X,c),u(X,c)]$。若 $\min_i \ell(X,c_i)\le U^2$（最小值精确计算），则该箱获证；
否则继续细分。区间无法判定的叶节点是\emph{未决}，必须如实报告，绝不能四舍五入成
「已覆盖」。历史构型最初的连续覆盖证书（certificate）即属此类：$R=153/1250$，
$22\,804$ 个叶节点，深度 $16$，$0$ 个未决，$21/21$ 项检查
（\texttt{verify/\allowbreak{}out/\allowbreak{}coverage\_\allowbreak{}n100.json}）。一个独立检查器从整数角点比较重新导出该判据，
与生成器不共享任何几何函数；它拒绝了三个刻意破坏的案例
（\texttt{verify/\allowbreak{}out/\allowbreak{}independent\_\allowbreak{}coverage\_\allowbreak{}n100.json}）。一个由 $12$ 个正例与负例组成的对照电池全部
通过（\texttt{verify/\allowbreak{}out/\allowbreak{}coverage\_\allowbreak{}controls\_\allowbreak{}n100.json}）。所有这些都是
\texttt{machine\_\allowbreak{}verified\_\allowbreak{}geometry = false}：从「箱已获证」到「$D$ 被覆盖」的归约是解析的。

\subsection{归约 B：格胞归约（精确，带有限枚举）}

对格构造可以做得比四叉树好得多：各胞铺满平面，因此 $D$ 的覆盖归约为对有限多个其胞
可能碰到 $D$ 的格点所做的有限精确检查。这就是 \S\ref{sec:lattice} 的方法，它产生的证书
其判定只用到整数／Fraction（有理数）。

\subsection{归约 C：外部构型的重放}

当构型由一个外部软件包给出时，诚实的上界路线是对该软件包自带检查器的逐字节忠实重放，
外加独立的交叉检查。这就是 \S\ref{sec:external} 的方法。

\section{三角格 $97{+}3$ 构型}
\label{sec:lattice}

\subsection{构型}

先在\emph{规范化}（normalized）坐标中工作，其中最近邻格点之间的距离为 $1$。格点是数对
$(i,j)\in\Z^2$，其实位置为
\[
x(i,j)=(u,\sqrt3\,v),\qquad u=i+\tfrac j2,\qquad v=\tfrac j2,
\qquad
\|x(i,j)\|^2=i^2+ij+j^2 .
\]
格点的六边形 Voronoi 胞是六边形 $H$ 的平移，其顶点在 $(u,v)$ 坐标下为
$(0,\pm\frac13)$、$(\pm\frac12,\pm\frac16)$，外接圆半径为 $1/\sqrt3$。

\begin{construction}[$97{+}3$ 构型]\label{con:lattice}
取所有满足 $i^2+ij+j^2<28$ 的格点（共 $97$ 个），再加上三个范数为 $28$ 的格点
\[
(-6,4),\quad(-2,-4),\quad(2,4).
\]
以缩放
\[
a=\sqrt{3/67},\qquad R=1/a=\sqrt{67/3},\qquad \alpha=1/\sqrt{67},
\]
以 $a\,x(i,j)$ 为缩放后圆心、半径为 $\alpha$ 的 $100$ 个闭圆盘覆盖 $D$，且 $\alpha$
\emph{恰好}是该圆心配置的覆盖半径（covering radius）。
\end{construction}

\begin{proposition}[构型的精确覆盖半径]\label{prop:lattice-alpha}
$\alpha=1/\sqrt{67}$ 是构造~\ref{con:lattice} 的 $100$ 个圆心覆盖 $D$ 的最小半径。
\end{proposition}

\begin{proof}[论证]
上半部分是下面描述的胞检查。下半部分考虑内部格洞 $q=(u,v)=(\tfrac12,\tfrac16)$，即被
六个格点围住的六边形中心。其规范化范数为
$\|q\|^2=u^2+3v^2=\tfrac14+\tfrac1{12}=\tfrac13$，而它到最近格点的距离就是六边形外接圆
半径 $1/\sqrt3$。按 $a$ 缩放后，最近保留下来的圆心距离为 $a/\sqrt3=1/\sqrt{67}=\alpha$，
而 $\|a q\|=a/\sqrt3=\alpha<1$，故 $aq$ 是 $D$ 中一点，它到每个圆心的距离恰好为
$\alpha$。因此没有更小的半径能覆盖 $D$。
\end{proof}

\begin{remark}
三个范数为 $28$ 的格点对覆盖是\emph{冗余}的：$97$ 个格点的配置已用同样的半径覆盖 $D$，
而证书记录了构造顺序（先 $97$ 个格点，再三个额外格点）。冗余并不能证明用 $100$ 个圆盘
可以达到更小的半径；它只说明这个特定构型没有用满它的预算。
\end{remark}

\subsection{精确的胞论证}

证书 \texttt{verify/\allowbreak{}out/\allowbreak{}lattice97\_\allowbreak{}upperbound\_\allowbreak{}n100.json}（模式
\texttt{n100-\allowbreak{}lattice97-\allowbreak{}upperbound-\allowbreak{}v1}）通过四个步骤把连续覆盖归约为有限精确检查，每一步都是整数或
\texttt{Fraction} 判定：

\begin{enumerate}[leftmargin=1.6em,itemsep=3pt]
\item \textbf{胞包含。} 格点 $p$ 的 Voronoi 胞中每一点都在 $p$ 的 $1/\sqrt3$ 之内
（$H$ 的外接圆半径）。
\item \textbf{有限格点表。} 若某格点的胞与半径 $R=\sqrt{67/3}$ 的闭圆盘相交，则该格点满足
$\|p\|^2<29$，从而 $\|p\|^2\le28$，从而 $|i|,|j|\le7$。所用枚举箱为 $|i|,|j|\le24$，
比充分箱 $7$ 宽裕得多。
\item \textbf{唯一顶点处的精确最小值。} 对每个满足 $\|p\|^2\le28$ 且 $p\notin C$ 的格点，
$\|x\|^2$ 在其胞上的最小值恰好为 $R^2=67/3$，并在唯一顶点 $v_p$ 处取到；这由支撑平面
不等式 $\langle v_p,\,w-v_p\rangle\ge0$（对胞的全部六个顶点 $w$）给出证书。
\item \textbf{外部壳层。} 满足 $\|p\|^2\ge29$ 的格点其胞距离
$\ge\sqrt{29}-1/\sqrt3>R$，由纯整数不等式 $3(29-22)^2>4\cdot29$，即 $7>2\sqrt{29/3}$
给出证书。
\end{enumerate}

\subsection{候选半径的代数数据}

\[
\alpha=\frac1{\sqrt{67}},\qquad P(t)=67t^2-1 \ \ (\text{coefficients } [67,0,-1]),
\]
$P$ 是本原的且首项系数为正；$67$ 不是平方数，故 $P$ 在 $\Q[t]$ 中不可约；并且
\[
\frac{1221694443}{10^{10}}<\alpha<\frac{1221694444}{10^{10}}
\]
是有理隔离区间。端点符号由整数平方比较判定，而 $P$ 在正半轴上严格递增
（$P'(t)=134t>0$），故该区间恰含一个根。这是\emph{本构型}覆盖半径的极小多项式——
它\emph{不}被声称是 $\rD(100)$ 的极小多项式，因为第三／第四部分的下界没有触及
$\alpha$。

\subsection{证书中的负对照}

格证书的十三条负对照必须全部被拒绝，也确实全部被拒绝了：一个放大的目标圆盘（仅作为
前提（precondition），明确\emph{不是}所交付命题的未覆盖见证（witness））；放大圆盘变体的
一个显式未覆盖点；一个缩小的覆盖半径；一个放大的覆盖半径；一个错误的度量；一个错误／
旋转过的六边形；一个膨胀的六边形；一次被截断的邻居搜索；一个归属错误的邻居；两个被删去的
相切圆心；把 $97$ 个圆心的构型当作「$100$ 点」命题提交；删去单个壳层圆心（覆盖仍然成立，
因此只有基数被拒绝）；以及一个被删去的内部圆心。一个独立实现
（{\footnotesize\texttt{verify/\allowbreak{}out/\allowbreak{}independent\_\allowbreak{}lattice97\_\allowbreak{}upperbound\_\allowbreak{}n100.json}}，\texttt{all\_\allowbreak{}passed = true}，
\texttt{coverage\_\allowbreak{}confirmed = true}）以不同的代码路径复现了该命题。

\section{外部 $100$ 圆心构型}
\label{sec:external}

当前被接受的最佳上界来自一个外部公开软件包，其提交被冻结在
\[
{\footnotesize{\footnotesize\texttt{67ebf1c5776f72a7b7786a8655eed5463429ae55}}},
\]
其字节已在本地原样重放，未作修改。它给出一个 $100$ 圆心的构型以及精确有理界
\[
U_\star=\frac{1154469548855829}{10^{16}}=0.1154469548855829 .
\]

\begin{center}
\footnotesize
\begin{tabular}{@{}>{\raggedright\arraybackslash}p{0.32\textwidth}>{\raggedright\arraybackslash}p{0.54\textwidth}@{}}
\toprule
产物 & 它所确立的内容 \\
\midrule
{\footnotesize\texttt{verify/\allowbreak{}out/\allowbreak{}external-\allowbreak{}disk100-\allowbreak{}bounds-\allowbreak{}67ebf1c.json}} &
完整的区间证据：退出码 $0$；六个下界分支（branch）全部重算为 \texttt{PASS} \\
{\footnotesize\texttt{verify/\allowbreak{}out/\allowbreak{}external-\allowbreak{}disk100-\allowbreak{}margin-\allowbreak{}67ebf1c.json}} &
一个独立的上界余量检查；退出码 $0$ \\
\texttt{verify/\allowbreak{}out/\allowbreak{}external-\allowbreak{}disk100-\allowbreak{}local-\allowbreak{}quadtree/\allowbreak{}replay.json} &
本项目自己的四叉树交叉检查，半径取略宽的 $R=0.115447$：主检查与独立检查均
\texttt{PASS}，$0$ 个未决 \\
\bottomrule
\end{tabular}
\end{center}

这里有两点评注。第一，同一个软件包还提供了第四部分所使用的\emph{六分支骨架}；吸收报告
固定了各自的角色：第三方材料是可重放的输入与外部参照，而不是被本项目自动继承的权威。
第二，重放是该软件包自带检查器的逐字节忠实运行，外加独立的交叉检查；它\emph{不}是
「该软件包的解析几何已被形式化」这一断言。

\section{协作搜索的收紧轨道}

在外部构型被吸收之前，本项目通过协作（自由点加退火）搜索收紧了它自己的构型。每个阶段
都把输出冻结为只读，并用独立的四叉树重放来认证一个有理上界：

\begin{center}
\footnotesize
\begin{tabular}{@{}>{\raggedright\arraybackslash}p{0.15\textwidth}>{\raggedright\arraybackslash}p{0.19\textwidth}>{\raggedright\arraybackslash}p{0.23\textwidth}>{\raggedright\arraybackslash}p{0.25\textwidth}@{}}
\toprule
阶段 & 浮点发现半径（不是判定） & 已认证的有理上界 & 证据 \\
\midrule
精确格 & $0.1221694443563052$ & $1/\sqrt{67}$（代数） &
\texttt{n100-\allowbreak{}algebraic-\allowbreak{}bounds-\allowbreak{}bundle.json} \\
协作，第 1 轮 & $0.1200703574813639$ & $18761/156250=0.1200704$ &
\texttt{n100-\allowbreak{}cooperative-\allowbreak{}anchor-\allowbreak{}tight/\allowbreak{}replay.json} \\
精化 & $0.1198908685151380$ & $119891/1000000=0.119891$ &
\texttt{n100-\allowbreak{}cooperative-\allowbreak{}refine-\allowbreak{}tight/\allowbreak{}replay.json} \\
抛光 & $0.1198764764337424$ & $119877/1000000=0.119877$ &
\texttt{n100-\allowbreak{}cooperative-\allowbreak{}polish-\allowbreak{}tight/\allowbreak{}replay.json} \\
\bottomrule
\end{tabular}
\end{center}

重放证书报告 \texttt{all\_\allowbreak{}passed = true}、\texttt{independent.verified = true}，抛光阶段有
$35\,584$ 个叶节点、最大深度 $23$、$0$ 个未决；冻结候选的十进制坐标被当作\emph{精确
有理数}，而不是替换为拟合出的格坐标或代数坐标。同一轨道的更早阶段认证了
$153/1250=0.1224$（阶段 1）与 $1223433/10000000=0.1223433$（阶段 2，服务器搜索，六个
候选，$15.746$ s 内 $916$ 次目标函数求值，六个候选都在连续 $150$ 次无改进之后停止）。

\subsection{协作轨道为何停止改进}

代数诊断（diagnostic）与界本身同样重要。每个阶段都恰好有一个紧洞（阈值 $10^{-9}$），
另有若干接近紧的洞参与一个有限带模型。抛光运行在触及 $4\times30$ 轮上限时仍记录了
$115$ 次严格改进，所以下降的终点尚未到达：该局部系统\emph{没有} KKT 或邻域最优性证书。
一次平稳性诊断运行了 $1.3$ s 并报告：一个 $10^{-6}$ 的近活跃带含 $123$ 个接触，但在
$10^{-7}$ 只有 $2$ 个接触、在 $10^{-9}$ 只有 $1$ 个；$114$ 行内部梯度，浮点秩 $114$，
零空间维数 $86$；并且在所测容差下找不到非负 KKT 乘子 LP。一个接触平衡／SQP 单元返回
\texttt{NO\_\allowbreak{}IMPROVEMENT}：它在 $60$ s 内把半径降低了约 $1.25\times10^{-9}$，低于接受
阈值，因此没有保存新候选，也没有更新上界。光滑极大方向大多增大了真实半径，被真实求值器
拒绝。所有这些都是有限浮点诊断：它既不证明不存在其他平衡分支，也不证明该候选
「不可代数化」。

\section{局部代数路线：一个孤立的 KKT 根}

独立于构型搜索，被吸收候选的局部代数也被重放了。孤立的代数 KKT 根具有

\begin{center}
\footnotesize
\begin{tabular}{@{}ll@{}}
\toprule
量 & 值 \\
\midrule
变量 & $932$ \\
基本变量 & $493$ \\
约束 & $439$ = $406$ 距离 + $32$ 边界 + $1$ 规范 \\
单项式 & $8\,348$，次数至多 $2$ \\
\bottomrule
\end{tabular}
\end{center}

证书 {\footnotesize\texttt{verify/\allowbreak{}out/\allowbreak{}external-\allowbreak{}disk100-\allowbreak{}strict-\allowbreak{}local-\allowbreak{}67ebf1c.json}} 记录了一个新根、
一次新的覆盖检查，以及 $8/8$ 个破坏测试被拒绝；一个 Krawczyk 式箱隔离该根，旋转商也被
处理。对同一根箱（半径 $10^{-58}$）的应力传递分析发现：箱内 $406$ 个距离乘子严格为正，
$32$ 个边界乘子严格为负，以及一个 $246$ 节点的连通接触图，其加权 Laplacian 块在 $214$
个自由节点上、以 $32$ 个边界接触为锚，具有精确的正定下界。然而，所尝试的 cover-19 式
全局 Dirichlet 能量下界失败了：由边界角余弦区间与杆图跳跃距离构造的 Loewner 矩阵仍产生
两个负的 LDL 主元，故该传递结果被记为 \texttt{H0\_\allowbreak{}PSD = false}。这个失败是一条路线
诊断，而不是对局部二阶条件的反例，也不是关于其他拓扑的陈述。

\section{上界\emph{不}给出什么}

\begin{itemize}[leftmargin=1.4em,itemsep=2pt]
\item 它们没有确定 $\rD(100)$：所接受的区间在绝对意义下宽度为 $0.0022$。
\item 它们没有证明格构型、协作候选或外部构型是最优的或局部最优的（唯一的例外是一个
孤立 KKT 根，那里认证了固定接触图的\emph{严格局部}极小性——一个关于单个图的局部陈述）。
\item 它们没有给出 $\rD(100)$ 的极小多项式。存在的是某个显式构型覆盖半径的极小多项式
$67t^2-1$，那是另一个对象。
\item 三个冗余的格圆盘表明这个构型没有用满其预算；这只是一个提示，而不是「用 $100$ 个
圆盘可能达到更小半径」的证明。
\end{itemize}

\part*{第三部分：普适解析下界（不依赖构型的路线）}
\addcontentsline{toc}{part}{第三部分：普适解析下界（不依赖构型的路线）}
\stepcounter{partanchor}\setcounter{section}{0}

\section{两类下界}

对 $\rD(100)$ 的下界可以用两种本质不同的方式得到。

\begin{enumerate}[leftmargin=1.6em,itemsep=3pt]
\item \textbf{普适（不依赖构型）界。} 一个单一的解析不等式，对某个半径范围内的任意
$100$ 个圆盘覆盖都成立，从而与低于某阈值 $L$ 的覆盖矛盾。特定构型的组合结构完全不
进入；代价是这类界相对较弱。
\item \textbf{候选半径矛盾。} 固定一个有理数 $R$ 并假设半径为 $R$ 的覆盖存在；随后导出
对所有可能构型的有限、穷尽的组合分类，并用精确不等式逐类排除。这是第四部分的引擎，
也是当前承载界 $0.113235$ 的东西。
\end{enumerate}

第三部分是第一种。它重要有三点原因：它给出一个不需要组合目录的地板；它校准了每个成分
至多能弥合多少差距；它产生了两个有教益的反驳（周界缺陷律（perimeter-defect law），
以及简单的全局能量），把纲领重新引向第四部分。

\section{面积界及其直接锐化}

最粗糙的界既精确又平凡：
\[
\pi=\area(D)\le\sum_{i=1}^{100}\area\bigl(B(c_i,r)\cap D\bigr)\le100\pi r^2
\quad\Longrightarrow\quad r\ge\frac1{10}.
\]
渐近密度陈述 $\delta=2\pi/(3\sqrt3)=1.2091996\ldots$ 在极限 $\eps\to0$ 下把 $1$ 换成
$\delta$，对边界可忽略的假想大 $n$ 构型给出 $\sqrt{\delta/100}=0.109964$。两个数都远低于
真值，而它们与所接受上界 $0.11544\ldots$ 之间的差距，就是每条严肃路线都必须挣得的
「边界修正」。

\section{路线：整数边计数，被修复的 Euler 路线}
\label{sec:integer-edge}

这是本纲领中写得最完整的解析路线。它证明精确陈述
\[
\rD(100)\ \ge\ L_1:=\frac{554570693}{5000000000}=0.1109141386 ,
\]
其意义在于：它修复而不是丢弃了历史上的「Euler + 边界胞（boundary cell）」链条——该链条
的审计曾把它降级为\emph{未证明}，并列出十条明确义务 O1--O10。

\subsection{建立：Voronoi 胞、弧与内接多边形}

假设 $100$ 个圆盘 $B(c_i,r)$ 以 $r<1$ 覆盖 $D$（若 $r\ge1$ 结论平凡），并施加第一部分的
投影与去重规范化（normalized），使得 $\|c_i\|\le1$ 且各 $c_i$ 互不相同。定义闭凸
Voronoi 胞
\[
V_i=\{x\in D:\|x-c_i\|\le\|x-c_j\|\ \ \forall j\}.
\]
三个事实会被反复使用：
\begin{itemize}[leftmargin=1.4em,itemsep=2pt]
\item $\bigcup_iV_i=D$，且对 $i\ne j$ 有
$\operatorname{int}V_i\cap\operatorname{int}V_j=\emptyset$；
\item $V_i\subseteq B(c_i,r)$：对 $x\in V_i$，由覆盖性有
$\|x-c_i\|=\min_j\|x-c_j\|\le r$；
\item 投影与去重之后每个 $V_i$ 都有正面积，因此不存在「空胞」记账问题。
\end{itemize}
设\emph{边界弧}（boundary arc）$A_1,\dots,A_{B_{\mathrm{arc}}}$ 为 $\partial D$ 的极大
正测度连通段，在其上归属（取得最小距离的最小下标）几乎处处为常数；孤立的零测度并列不
切断弧。设\emph{边界顶点}为各弧的端点，并设
\[
P:=\conv(\text{all boundary vertices})
\]
为内接多边形。

\begin{lemma}[弧长界；以及一处勘误]\label{lem:arc}
对 $r<1$，每条弧的角测度 $\le2\arcsin r<\pi$。从而
$B_{\mathrm{arc}}\ge\pi/\arcsin r$，且没有弧能绕满整个圆周。
\end{lemma}
\begin{proof}
若某条弧的测度 $\psi\ge\pi$，它就含一条测度为 $\pi$ 的子弧，其两个端点对径，故距离为
$2$；两者都落在 $V_i\subseteq B(c_i,r)$ 中，于是 $2r\ge2$，即 $r\ge1$，矛盾。对
$\psi<\pi$ 弦长公式适用：弦长 $\ell\le2r$ 时 $\psi=2\arcsin(\ell/2)\le2\arcsin r<\pi$。
对各弧求和得 $2\pi\le B_{\mathrm{arc}}\cdot2\arcsin r$。
\end{proof}

\begin{remark}[一个被更正并留档的常数]
早期草稿用 $4\arcsin(r/2)$ 作为 $\partial D\cap B(c_i,r)$ 角测度的上界。那是错的：它只是
$\|c_i\|=1$ 处的取值，\emph{小于}真正的最大值。令
$Q(\rho)=(1+\rho^2-r^2)/(2\rho)$，测度为 $2\arccos Q(\rho)$，而 $Q$ 在
$\rho_\ast=\sqrt{1-r^2}$ 处取极小，给出最大值
$2\arccos\sqrt{1-r^2}=2\arcsin r$。错误常数无法对接近 $1$ 的 $r$ 推出 $<\pi$，因此不可
用；勘误保留在源文件中。
\end{remark}

\begin{lemma}[引理 A：圆盘内凸多边形的面积，圆心不受限]\label{lem:A}
设 $K\subset\R^2$ 凸且紧，$K\subseteq B(c,r)$，其中 $c$ 不必属于 $K$；设 $k$ 为 $K$ 的
极点个数。则
\[
\area(K)\le r^2f(k),\qquad
f(k)=\tfrac k2\sin\tfrac{2\pi}k\ \ (k\ge3),\qquad f(0)=f(1)=f(2):=0 .
\]
\end{lemma}
\begin{proof}[五步梗概]
设 $B(o,R)$ 为 $K$ 的最小包围圆，故 $R\le r$。接触集 $S=K\cap\partial B(o,R)$ 由顶点
组成且非空。一个有限分离论证表明 $o\in\conv(S)$：否则把 $o$ 朝 $\conv(S)$ 略微平移，会
严格减小到每个接触顶点的距离，并由有限性也减小到其他每个顶点的距离，与极小性矛盾。故
$o\in K$。由于 $o\in K$，以极点 $v_1,\dots,v_m$ 为顶点的三角形 $(o,v_j,v_{j+1})$ 恰好
剖分 $K$（其中边界与顶点位置的 $o$ 另行处理；「$o$ 是极点」的情形被排除，因为在开度
$<\pi$ 的凸锥内不可能有非负组合 $\sum_s\lambda_su_s=0$）。因此
$\area(K)\le\frac{R^2}2\sum_j\sin\alpha_j$，其中 $\alpha_j\in[0,\pi]$ 且
$\sum_j\alpha_j=2\pi$。对凹的正弦函数用 Jensen 不等式得
$\sum_j\sin\alpha_j\le m\sin(2\pi/m)$，而 $M\mapsto M\sin(2\pi/M)$ 对 $M\ge3$ 递增，
因为 $x\mapsto\sin x/x$ 在 $(0,\pi)$ 上严格递减。最后 $m\le k$，而 $f$ 在整数 $\ge3$ 上
递增；$m\le2$ 时面积为 $0$。
\end{proof}

\begin{remark}
引理~\ref{lem:A} 去掉了历史假设 $c_i\in V_i\cap P$。扇形分解以最小包围圆的圆心 $o$ 为
基准，而 $o\in K$ 是被\emph{证明}的，不是被假设的。该界是最优的：对半径为 $r$ 的圆的
内接正 $k$ 边形取等号。
\end{remark}

\begin{lemma}[引理 B：帽属于它们自己的胞；$Q_i=V_i\cap P$ 精确成立]\label{lem:B}
对每条弧 $A$，令 $C(A)=\conv(A)$ 为其弦所截的弓形，并令 $H_A^{\mathrm{in}}$ 为以该弦为
边界、包含原点的闭半平面。
\begin{enumerate}[leftmargin=1.6em,itemsep=2pt]
\item 若弧 $A$ 属于胞 $V_i$（即 $A\subseteq V_i$），则 $C(A)\subseteq V_i$。
\item 对 $i\ne j$ 以及胞 $j$ 的弧 $A$，有 $V_i\subseteq H_A^{\mathrm{in}}$。
\item 从而，令 $Q_i:=V_i\cap P$，
\[
Q_i=V_i\cap\bigcap_{A\ \text{of cell}\ i}H_A^{\mathrm{in}}
\quad\text{(exact identity)},
\]
\[
\bigcup_iQ_i=P,\quad \operatorname{int}Q_i\ \text{pairwise disjoint},\quad Q_i\subseteq B(c_i,r).
\]
\item $\#\text{sides}(Q_i)\le m_i+a_i=:e_i$，其中 $m_i$ 是 $V_i$ 的极大直边数，$a_i$ 是
胞 $i$ 所拥有的弧数；且 $\sum_ia_i=B_{\mathrm{arc}}$。
\end{enumerate}
\end{lemma}
\begin{proof}[梗概]
(1) 由两行凸性论证得 $C(A)=\conv(A)$，而 $V_i$ 凸。(2) 若
$x\in V_i\cap\operatorname{int}C(A)$，则因
$\operatorname{int}C(A)\subseteq
\operatorname{int}V_j$ 且不同胞的内部互不相交，得到
矛盾；边界情形沿从 $V_i$ 内点出发的一条线段取极限即得。(3) 取所有半平面的交，并由 (2)
删去冗余者。(4) 每条被拥有的弦半平面至多增加一条直边，且不能切开已有的直边；退化的
$Q_i$（点或线段）面积为 $0$。
\end{proof}

\begin{remark}[风险点，已作答]
自然的担心是 $V_i\cap P$ 可能获得来自\emph{其他}胞之弦的边。引理~\ref{lem:B}(2) 表明
不会：每个帽都包含在其自身胞的凸包内，所以 $V_i\cap P$ 与 $V_i$ 的差别恰好是它自己的
帽。这对任意弧数、任意退化情形与任意 $r$ 都成立。
\end{remark}

\begin{lemma}[引理 C：$P$ 的面积下界]\label{lem:C}
若 $g_1,\dots,g_B\ge0$，$\sum_kg_k=2\pi$ 且 $g_k\le\gamma\in(0,\pi]$，则
\[
A:=\tfrac12\sum_k\sin g_k\ \ge\ \pi\frac{\sin\gamma}{\gamma}\ \ge\ \pi\Bigl(1-\frac{\gamma^2}6\Bigr)
\qquad(\gamma<1).
\]
\end{lemma}
\begin{proof}
$\sin g_k=g_k\,s(g_k)\ge g_k\,s(\gamma)$，因为 $s(t)=\sin t/t$ 在 $(0,\pi]$ 上递减；
求和即得。第二个不等式是 $\sin\gamma/\gamma$ 的交替级数界。
\end{proof}

\begin{lemma}[引理 D：平面图边记账（edge accounting），无需一般位置]\label{lem:D}
构造边界 Voronoi 复形：顶点为内部 Voronoi 顶点（$\ge3$ 个胞以正楔角相交处）与
$B_{\mathrm{arc}}$ 个边界顶点；边为\emph{在每个图顶点处切断}的平分线段，加上
$B_{\mathrm{arc}}$ 条弧边；面为 $N$ 个胞加上外部面。则每个顶点的度 $\ge3$，且
\[
\sum_{i=1}^{N}e_i\ \le\ 6N-6-B_{\mathrm{arc}} .
\]
\end{lemma}
\begin{proof}[梗概]
度数：在内部顶点处，相交的胞填满整个 $2\pi$ 的楔角，这些楔角被正长度的平分射线隔开，
故至少三个；在边界顶点处，来自两侧的弧端点加上至少一条进入 $\operatorname{int}D$ 的
平分射线给出度 $\ge3$（平分线与 $\partial D$ 横截相交，因为
$\varphi(\theta)=\|u(\theta)-c_k\|^2-\|u(\theta)-c_i\|^2$ 的符号在此改变，故
$\varphi'\ne0$）。不假设连通性的 Euler 公式，在 $K\ge1$ 个连通分支时读作
$E_s=V_{\mathrm{int}}+N-K\le V_{\mathrm{int}}+N-1$。双重计数给出
$\sum_im_i\le2E_s$，且 $\sum_ia_i=B_{\mathrm{arc}}$。于是由
$2E_s\ge3V_{\mathrm{int}}+B_{\mathrm{arc}}$ 得
$V_{\mathrm{int}}\le2N-2K-B_{\mathrm{arc}}$，
$2E_s\le6N-6K-2B_{\mathrm{arc}}\le6N-6-2B_{\mathrm{arc}}$，最终
$\sum_ie_i\le6N-6-B_{\mathrm{arc}}$。
\end{proof}

\begin{remark}[表明假设不可去掉的精确反例]
(i) $n=1$，$c_1=0$，$r=1$：单条弧绕满圆周，两个端点重合，恒等式
$\#\text{boundary vertices}=B_{\mathrm{arc}}$ 失效，且 $\sum m_i=0$ 违反
$6N-6-2B_{\mathrm{arc}}=-2<0$。由引理~\ref{lem:arc}，这种退化对 $r<1$ 不可能。
(ii) 三圆心构型 $c_1=0$，$c_2=(a,0)$，$c_3=(-a,0)$（$0<a<2$）中
$V_1=D\cap\{|x_1|\le a/2\}$ 拥有\emph{两条}弧，$V_2,V_3$ 是仅有帽的胞且 $e=2$，故
$B_{\mathrm{arc}}=4$，而触及边界的胞数为 $3$；且
$\sum_ie_i=8=6N-6-B_{\mathrm{arc}}$ 取到等号。把 $B$ 读作边界胞数会损失一个单位的
强度；假设「每个边界胞一条弧」会把 $e_1$ 记为 $3$ 而不是 $4$，这根本是错的。
\end{remark}

\subsection{七步聚合}

假设存在半径 $r\le L_1$ 的覆盖，且 $r<1$。

\begin{enumerate}[leftmargin=1.8em,itemsep=3pt]
\item[(1)] \emph{边界结构。} 由引理~\ref{lem:arc}，
$B_{\mathrm{arc}}\ge\pi/\arcsin r$；因 $\pi/\arcsin L_1=28.2663>28$ 且
$B_{\mathrm{arc}}$ 为整数，$B_{\mathrm{arc}}\ge29$。又 $B_{\mathrm{arc}}\ge3$，故 $P$
是真正的多边形。
\item[(2)] \emph{剖分与面积恒等式。} 由引理~\ref{lem:B}，
$\area(P)=\sum_i\area(Q_i)$，其中 $N\le100$，每个 $Q_i\subseteq B(c_i,r)$ 凸且
$\#\text{sides}(Q_i)\le e_i$。
\item[(3)] \emph{逐胞面积界。} 由引理~\ref{lem:A} 以及 $f$ 在整数 $\ge3$ 上的单调性，
$\area(Q_i)\le r^2f(e_i)$；$k_i\le2$ 的胞贡献 $0=r^2f(e_i)$。故
$\area(P)\le r^2\sum_if(e_i)$。
\item[(4)] \emph{边记账。} 引理~\ref{lem:D} 给出
$\sum_ie_i\le6N-6-B_{\mathrm{arc}}\le600-6-29=565$。
\item[(5)] \emph{仿射割线记账。} 令
\[
a:=f(6)-f(5)=0.2204349206\ldots>0,\qquad
b:=6f(5)-5f(6)=1.2754666876\ldots>0,
\]
并令 $\varphi(e)=ae+b$ 为过 $(5,f(5))$ 与 $(6,f(6))$ 的割线。对整数 $e\ge3$ 有
$f(e)\le\varphi(e)$：$e\ge6$ 时这来自 $f$ 在 $[3,\infty)$ 上的凹性
（$f''(M)=-2\pi^2\sin(2\pi/M)/M^3<0$），而 $e=3,4$ 可精确验证，
\[
\varphi(3)-f(3)=\tfrac{15}4\bigl(2\sin72^\circ-\sqrt3\bigr)>0,\qquad
\varphi(4)-f(4)=5\sin72^\circ-\tfrac32\sqrt3-2=0.1568\ldots>0 .
\]
$e\le2$ 的胞被删除（它们贡献 $0$）。故
\[
\sum_if(e_i)\le a\sum_ie_i+bN\le565a+100b=65f(6)+35f(5),
\]
最后一步是代数恒等式 $565(f_6-f_5)+100(6f_5-5f_6)=65f_6+35f_5$。数值上
$565a+100b=252.0923989138\ldots$。这一手段把「整数凹性加上 $92$ 次单独符号检查」替换为
两个有限不等式，而且关键在于它是\emph{严格的}：连续松弛会给出
$100f(5.65)=253.2941376702\ldots$，所以离散记账买到了 $1.2017387564$
（$\approx0.48\%$）。
\item[(6)] \emph{面积下界。} 由引理~\ref{lem:C} 与~\ref{lem:arc}，取
$\gamma=2\arcsin r$，$\area(P)\ge\pi(1-\gamma^2/6)$。
\item[(7)] \emph{收尾。} 把 (3)--(5) 与 (6) 结合，
\[
\pi\bigl(1-\tfrac16(2\arcsin r)^2\bigr)\ \le\ r^2\bigl(65f(6)+35f(5)\bigr)
\qquad\text{for all }r\le L_1 .
\]
左端关于 $r$ 递减而右端关于 $r$ 递增，故只需在 $r=L_1$ 处反驳该不等式。精确地，
\[
65f(6)+35f(5)=\frac{195\sqrt3}{2}+\frac{175}{8}\sqrt{10+2\sqrt5},
\]
\[
\pi\bigl(1-\tfrac16(2\arcsin L_1)^2\bigr)=3.1157211653\ldots
\ >\ L_1^2\cdot252.0923989138\ldots=3.1012271140\ldots,
\]
余量为 $0.0144940513$（相对意义下 $0.467\%$）。最后一步是一个有限的代数／有理不等式，
由算术代理用 $\pi$ 与 $\arcsin L_1$ 的有理包络做了 $15$ 次精确检查而闭合。因此不存在
半径 $\le L_1$ 的覆盖：
\[
\rD(100)\ \ge\ L_1=\frac{554570693}{5000000000}=0.1109141386 .
\]
\end{enumerate}

\begin{remark}[修复的内容]
历史链条使用了 $\area(P)_{\min}(r,B)$ 的精确极值以及一个连续不动点，其根位于 $L_1$ 下方
$7.6\times10^{-12}$ 处。修复后的路线刻意使用\emph{更粗糙}的下界 $\pi(1-\gamma^2/6)$，
放弃约 $0.46\%$ 的精度，换来的是一举消除关于 $B$ 的取值以及 $P$ 的顶点分布的所有
问题——正是这一点使最终不等式既严格又有限。$0.467\%$ 的净余量完全来自第 (5) 步严格的
离散割线记账。
\end{remark}

\begin{center}
\footnotesize
\begin{tabular}{@{}>{\raggedright\arraybackslash}p{0.28\textwidth}>{\raggedright\arraybackslash}p{0.60\textwidth}@{}}
\toprule
历史义务 & 修复后的状态 \\
\midrule
O1 退化的度 $\ge4$ 顶点 & 已吸收：只需度 $\ge3$ 加上不等式
$\sum_im_i\le2E_s$ \\
O2 退化的边界顶点改变 $b$ & $B$ 定义为弧数；对 $r<1$ 每条弧
$<\!\pi$ \\
O3 $b$ 的良定义性、多弧胞 & 已解决：$e_i=m_i+a_i$，并有一个精确反例 \\
O4 扇形分解需要 $c_i\in V_i\cap P$ & 已去除：引理~\ref{lem:A} 使用最小包围圆
圆心 \\
O4$'$ 帽项界 & 已去除：帽归其所在胞所有（引理~\ref{lem:B}） \\
O5 圆心在 $D$ 外 & 投影规范化；不再需要 $c_i\in P$ \\
O6 圆心重合 & 去重；每个胞都有正面积 \\
O7 共线三元组／平行平分线 & 已解决：只需 Euler、度 $\ge3$ 与 $\le$ 方向 \\
O8 冗余圆盘／空胞 & 投影与去重后没有空胞 \\
O9 边界退化（相切、单点弧） & 零测度接触不产生弧；绕满圆周的弧被排除 \\
O10 不动点／单调性升级 & 已去除：替换为 $r=L_1$ 处的一个有限不等式 \\
\bottomrule
\end{tabular}
\end{center}

\subsection{这条路线中仍然开放的部分}

几何部分依赖于一个有限不等式，其精确算术被委派出去；而「在边界顶点相交的胞填满半个
邻域」这一步是没有机器证书的几何论证。该路线不假设一般位置，因此不需要扰动或极限论证；
若仍要使用，正确的形式是取 $\eps<(L-r)/2$ 并把半径放大到 $r+\eps<L$，\emph{同时}由
$\eps$ 的选取重新保证 $r+\eps\le L$——要避免的错误是扰动越过 $L$ 却继续使用
$r+\eps\le L$。

\section{路线：锐化包络（sharpened envelope）}

同一骨架——投影、凸 Voronoi 胞、多弧 Euler 记账、整数边计数——随后用重新认证的参数推进
\[
\arcsin L<T,\qquad 28T<\pi,
\]
它们强制 $B\ge29$ 与 $\sum_ie_i\le565$，而面积上界系数保持其精确形式
$G=65f_6+35f_5$，面积下界保持更锐的 $\pi\sin(2\arcsin r)/(2\arcsin r)$（而不是截断
级数）。一个交替有理多项式下界随即一次性排除所有 $0<r\le L$。结果是
\[
\rD(100)\ \ge\ L_2:=\frac{1111720153}{10^{10}}=0.1111720153 ,
\]
比前一个值恰好改进 $2578767/10^{10}=0.0002578767$，其正余量约
$2.5235222738557335\times10^{-9}$，以精确分数存于证书中。主算术执行 $16$ 次精确检查；
独立算术使用 $\arctan$ 与平方根包络，与主算术不共享任何级数实现，并通过 $40$ 次检查加
$9$ 个负对照。该路线的理想包络有唯一零点，隔离于
$(1111720153/10^{10},\,1111720154/10^{10})$——这是关于\emph{方法极限}的陈述，而不是
$\rD(100)$ 的区间。更高精度无法沿这条路线弥合与上界的差距。来源：
\texttt{proofs/\allowbreak{}n100-\allowbreak{}sharp-\allowbreak{}envelope-\allowbreak{}lowerbound.md}、
\texttt{docs/\allowbreak{}review-\allowbreak{}sharp-\allowbreak{}envelope-\allowbreak{}n100.md}、
\texttt{verify/\allowbreak{}out/\allowbreak{}n100-\allowbreak{}algebraic-\allowbreak{}bounds-\allowbreak{}bundle.json}。

\section{路线：长弦二次亏损（long-chord quadratic deficit）}

另一个普适界来自与边界弧的长弦相关的二次亏损。经过三轮独立评审后，几何草稿被冻结于
SHA-256
{\footnotesize{\footnotesize\texttt{658863d2a36118be915fda66f10596f05073aadd51738da55bcf35f967d8cf73}}}，
主检查器与独立有理算术检查器通过
\texttt{verify/\allowbreak{}out/\allowbreak{}n100-\allowbreak{}longchord-\allowbreak{}bundle-\allowbreak{}frozen.json} 重放：统一退出码 $0$，子退出码 $[0,0]$，
\texttt{all\_\allowbreak{}passed = true}，\texttt{universal\_\allowbreak{}lower\_\allowbreak{}bound\_\allowbreak{}accepted = true}，全部 $566$ 个弧数
被枚举且余量为正，最坏情形 $B=42$，$6/6$ 个负对照被拒绝。这给出
\[
\rD(100)\ \ge\ L_3:=\frac{2787}{25000}=0.11148 ,
\]
它同时优于 $L_1$ 与 $L_2$，是当前被接受的最强免构型界。其几何是解析证明加评审
（\texttt{machine\_\allowbreak{}verified\_\allowbreak{}geometry = false}）；它不是形式化证明。

\begin{remark}[冻结包内部的分层，保留可见]
同一个包在其自身的子块中比在顶层更保守，这一区别很重要。顶层字段报告
\texttt{universal\_\allowbreak{}lower\_\allowbreak{}bound\_\allowbreak{}accepted = true}，并且主算术行与独立算术行在 $566$ 行上一致给出约
$4.4799\times10^{-4}$ 的最小余量。然而在几何子块内部，
\texttt{all\_\allowbreak{}geometric\_\allowbreak{}premises\_\allowbreak{}closed = true}，而
\texttt{combined\_\allowbreak{}universal\_\allowbreak{}lower\_\allowbreak{}bound\_\allowbreak{}accepted = false}，并且算术被标记为
\texttt{reported\_\allowbreak{}exit\_\allowbreak{}0\_\allowbreak{}by\_\allowbreak{}parent\_\allowbreak{}agent\_\allowbreak{}not\_\allowbreak{}recomputed\_\allowbreak{}and\_\allowbreak{}not\_\allowbreak{}accepted\_\allowbreak{}here}：对一个数值阈值
断言而言，该链条还需要前提 (P1) 与 (P2) 由精确有理算术\emph{以及}一次独立检查来认证，而
这额外的一步位于单独的算术证书中，而不是在冻结的几何块内部。引用顶层的接受而不加这个
限定，会夸大几何块自身所认证的内容。
\end{remark}

\section{被反驳的普适路线}

\subsection{周界缺陷（$\lambda$）律不可能成立}

早期纲领计划提出证明如下形式的律
\[
n r^2\ \ge\ \delta\,(1+\lambda r),\qquad \delta=\frac{2\pi}{3\sqrt3}=1.2091996\ldots,
\]
并调节 $\lambda$，使 $n=100$ 被钉在 $0.4\%$ 以内。该计划已死，原因是一个两行计算。若
这样的律对某个 $n$ 成立，则在该 $n$ 已被证明的最优值 $r_n$ 处它强制
\[
\lambda\ \le\ \frac{n\,r_n^2/\delta-1}{r_n},
\]
而被证明的取值给出 $0.697$（$n=8$）、$\mathbf{0.669}$（$n=9$）、$0.734$（$n=10$）、
$0.752$（$n=19$）、$0.763$（$n=25$）。所以 $\lambda\approx0.87$ 的全局律不可能成立——
仅 $n=9$ 就反驳了它——而最大可容许常数 $\lambda^\ast=0.6688$ 至多给出
$\rD(100)\ge0.114081$，即比已知最佳构型低 $6.75\%$。更糟的是，本应提供周界层的
面积--周界实现从未逐胞起约束作用（$0/25$ 个案例），其聚合上界仅为 $0.1540$。原始校准表
（显示值）为：$\lambda=0.76$ 时 $0.114655$，$\lambda=0.83$ 时 $0.115096$，
$\lambda=0.86$ 时 $0.115286$——它们全都不可达，因为 $n=9$ 的天花板是 $0.669$。

该路线在 $n=25$ 上校准后被正式降级：同一方法在那里以有效普适界 $0.2235441655$ 闭合
（即距精确值 $-6.52\%$，或按最有利的读法 $-5.75\%$）；在 $n=100$ 上同一路线为
$-9.34\%$，其周界缺陷常数为 $\lambda\approx0.15$，而阈值为 $\approx1.01$。那里用到的
三个引理是边界 Euler 记账 $\sum(6-m_i)=6+2b$、「边界胞 = 一条额外边」
$\max_\Theta g=f(m+1)$，以及弧覆盖界 $b\ge\pi/\arcsin r$；其中第一个正是后来被
引理~\ref{lem:D} 用更安全的不等式 $\sum_ie_i\le6N-6-B_{\mathrm{arc}}$ 替换掉的东西。
换言之，\S\ref{sec:integer-edge} 的现代普适界，是一条以原形被放弃的路线的修复后裔。

\subsection{见证（witness）校准：正确但太弱}

一条有限见证路线被完整认证，随后搁置。它有 $21$ 个有理见证点、$32$ 个极大模式、$53$ 个
可行非空子集与 $506$ 次枚举事件，总权重 $21/2$，在 $r=7/25$ 处排除 $n=8,9,10$。独立
检查器通过 $74$ 次检查，$0$ 次失败。但该界弱于 $n=8,9,10$ 的已知面积界，也没有改进
$n=100$ 的区间：在这个网格上 $r=7/25$ 处的对偶最优值恰好为 $12$，而对 $r\ge3/10$，
八个圆盘已经覆盖了见证集。因此该构造无法越过面积界成为一条路线，而有限归约的解析完备性
依赖于一个未形式化的平面 Helly／最小包围圆支撑界。

\subsection{简单的全局能量在候选图上失败}

曾提出一条角／调和能量路线，把被吸收的 $n=19$ 证明的应力结构搬到 $n=100$。两份证书
了结此事。第一，真正的调和延拓应力能量的二次型，在候选角域内的候选根处其区间上界严格
低于 $-1/100000$
（{\footnotesize\texttt{verify/\allowbreak{}out/\allowbreak{}n100-\allowbreak{}angular-\allowbreak{}energy-\allowbreak{}independent.json}}，退出码 $0$）。第二，一个有理
旋转反例表明，在允许的角域内能量可以\emph{减少}超过 $10^{-12}$
（{\footnotesize\texttt{verify/\allowbreak{}out/\allowbreak{}n100-\allowbreak{}harmonic-\allowbreak{}energy-\allowbreak{}relaxation-\allowbreak{}counterexample.json}}，退出码 $0$）。
因此固定应力的全局凸能量路线不可用，必须替换为逐杆的力与约束分支（branch），并配以
定量隔离
（{\footnotesize\texttt{verify/\allowbreak{}out/\allowbreak{}n100-\allowbreak{}angular-\allowbreak{}energy-\allowbreak{}obstruction-\allowbreak{}audit.json}}）。这是关于一个候选
接触图的陈述；它不是覆盖反例，也不与局部二阶条件矛盾。

\subsection{另外两条未被提升的数值路线}

\begin{itemize}[leftmargin=1.4em,itemsep=2pt]
\item 一个逐胞等弧／边数-$5,6$ 抽样族找到 $0.1118772$，但对手族的完备性从未被证明，
因此它不能作为普适界。
\item 一个径向壳层 LP 报告了接近 $92.56$ 的值；细网格复核发现所抽样的解不可行，所以它
不是覆盖数的已认证界。该报告在独立质疑下被更正。这并不证明有限见证路线不可能——只说明
所测试的松弛没有兑现。
\end{itemize}

\part*{第四部分：六分支引擎，固定半径下的矛盾}
\addcontentsline{toc}{part}{第四部分：六分支引擎，固定半径下的矛盾}
\stepcounter{partanchor}\setcounter{section}{0}

\section{逻辑骨架}

本部分的全部内容都具有如下形态。

\begin{enumerate}[leftmargin=1.6em,itemsep=3pt]
\item 固定一个有理的\emph{候选半径} $R$，并反设 $100$ 个半径为 $R$ 的闭圆盘覆盖 $D$。
\item \textbf{规范化（normalization）}该半径下的每个圆心（同半径规范化），并重建归属关系、有效使用的圆盘数与公共三角剖分。先前半径的任何东西都不得复用。
\item \textbf{计数。}导出必要条件，排除所有至多使用 $m-1$ 个有效圆盘的构型；并对 $m=100$ 把边界弧（boundary arc）数 $B$ 限制为一张有限清单。
\item \textbf{分裂为分支（branch）。}存活的 $(m,B)$ 对划分为六个互斥分支 $B31$、$B32$、$B33$、$B34$、$B35$、$B{\ge}36$。
\item \textbf{闭合每个分支}：用一个面积/曲率预算，即一条与覆盖性相矛盾的精确有理不等式。
\item \textbf{合并。}合并器检查：每个分支都在同一半径下闭合；预算只被花费一次；任何在\emph{更大}半径下闭合的分支都通过半径单调性（monotonicity inheritance）被继承。
\end{enumerate}

最后一点是一条定理，而非一项约定：

\begin{lemma}[半径单调性]\label{lem:monotone}
若不存在用 $100$ 个半径为 $R'$ 的圆盘作出的覆盖，则不存在用 $100$ 个半径为 $R<R'$ 的圆盘作出的覆盖。
\end{lemma}
\begin{proof}
半径 $R$ 下的覆盖在把每个圆盘放大之后即为半径 $R'$ 下的覆盖，因为 $B(c,R)\subseteq B(c,R')$。
\end{proof}

因此，在 $R'=2831/25000$ 处建立的分支证书（certificate）对任意 $R\le R'$ 都可用，前提是合并器显式记录该继承（D 级证据）。在当前证书链中，尾部分支\emph{不}以这种方式继承：它在 $R_\star=22647/200000$ 处被直接重放。

\section{固定半径下的规范化}

\subsection{径向映射}

对候选半径 $R$，定义径向映射
\[
d'=\frac{1-R^2}{d},
\]
其中 $d$ 是原点到某圆心的距离，并令 $u$ 为该圆心的单位向量。恒等式
\begin{equation}\label{eq:radial-identity}
\|z-d'u\|^2-R^2=\frac{d'}{d}\bigl(\|z-du\|^2-R^2\bigr)+\Bigl(1-\frac{d'}{d}\Bigr)\bigl(\|z\|^2-1\bigr)
\end{equation}
在证书中由 CAS 精确断言。由于 $d'/d\le1$，且对 $z\in D$ 有 $\|z\|^2-1\le0$，每个在半径 $R$ 下已被覆盖的点在该映射之后仍被覆盖。规范化在进入计数分支\emph{之前}施加于\emph{全部}圆心；归属关系、有效圆盘数与公共三角剖分都从头重建。这正是引擎无法通过修改一个常数就在新半径下复用的原因：图本身发生了变化。

\subsection{公共计数桥（counting bridge）}

在 $R=11323/100000$ 处，计数桥排除所有至多使用 $99$ 个有效圆盘的构型：$n=99$，$96$ 个情形，最紧的 $B=34$，其精确最小面积余量为
\[
\resizebox{\textwidth}{!}{$\displaystyle \frac{93591732347747920006153107439494216107661633777884476873367781737690463949606826459158012299304439}{17\cdot10^{100}}>0$}
\]
（存放在外部重放证书中）。对 $m=100$，存活的边界计数恰好是
\[
B\in\{31,32,33,34,35,36,37,38,39\}
\]
（\texttt{complete\_\allowbreak{}surviving\_\allowbreak{}boundary\_\allowbreak{}counts}）。六个分支划分这张清单，并且按 $B$ 的定义它们互斥。

同一条计数桥在 $R=2831/25000$ 处被重新导出：全部 $m\le99$ 都被排除，而对 $m=100$ 只有 $B31$--$39$ 存活。规范化的圆心范数平方上界，即 $1-R^2$ 的一个向外有理包围，是
\[
1-R^2\ \le\ \frac{9871789671}{10^{10}}\quad(R=11323/100000),\qquad
1-R^2\ =\ \frac{616985439}{625000000}\quad(R=2831/25000),
\]
而在被接受的半径 $R_\star=22647/200000$ 处有 $1-R_\star^2=0.987177834775\ldots$；被接受的证书对更小的半径 $1132301/10^7$ 使用有效的外向界 $98717894445399/10^{14}$。这些是同一个量在不同半径下的取值，每一个都与它自己的半径一起给出。

\subsection{分支词汇}

\begin{center}
\footnotesize
\begin{tabular}{@{}lp{0.66\textwidth}@{}}
\toprule
符号 & 含义 \\
\midrule
$B$ & 所服务的边界弧数（等价地，公共图的边界顶点数） \\
$m$ & 有效使用的圆盘数 \\
$k$ & 不同的边界服务圆盘数 \\
$M$ & 规范化的公共三角剖分的内部面数 \\
$e$、$b$ & 耳（ear）数，相应地坏面（bad face）数（「$e0$」「$b0$」表示为零） \\
$q$ & 一个顶点的标记面（marked face）数；$q1$ 即「坏 $q1$」 \\
$h$ & 严格深点（deep vertex）的总度超出量 \\
$C$ & 核心（core）边界的弦数 \\
$a$ & $C-f_3$，其中 $f_3$ 是空的核心边界三角形数 \\
DDD, DDR, DRR, BRR, RRR & 按其三个顶点的度（深点/环点（ring vertex））划分的面类型 \\
核心 $P/F$ & 诱导核心圆盘的边界顶点数 / 面数 \\
\bottomrule
\end{tabular}
\end{center}

每个分支中反复出现的机制是一个\emph{预算}：半径 $R$ 处可用的边界弧长必须为内部面积、边界「标记」三角形、耳与坏面亏空，以及一次铰链（hinge）预算或方差（variance）预算支付买单。全程强制两条结构规则：(i) 由两个面共享的一条物理边（physical edge）只被支付\emph{一次}，相反的半边项精确抵消；(ii) 一个标量亏空或一项局部增益不得在两行中重复花费。证书把这些记为显式字段（\texttt{physical\_\allowbreak{}edge\_\allowbreak{}pairs}、\texttt{hinge\_\allowbreak{}payment = 0}、\texttt{variance\_\allowbreak{}payment = 0}、budget-review 字符串）。

\section{分支 B31}

$B31$ 是最小的边界计数。它的证书是「全局全面曲率」行：$40$ 个面支撑与 $32$ 个顶点支撑、一个度窗口、就地收费、耳与坏面系数，以及两个 $M$-区间（$M\le23$ 与 $M\ge24$）。值得注意的结构字段：全部二十个面标签多重集都被允许，不假设任何深点标签（\texttt{deep\_\allowbreak{}label\_\allowbreak{}assumption = false}），不使用反射圆心假设，平方边罚与半侧乘子都为 $0$，并且允许任意的耳与坏面。该分支还带有一个 $n=99$ 排除，含 $96$ 个情形、最紧的 $B=34$，其余量等于计数桥的余量。

转移到新半径的几何骨架在证书本身中被逐条写为义务：同半径的整覆盖规范化与图重建；无条件的全星形度损失不等式及其预算比较；整圆盘 $/6$ 与诱导核心 $/5$ 的等周（isoperimetry）；相同的全面标签配以带符号的边抵消与一次性预算（one-shot budget）使用；以及一个闭端点有限型连续性论证。

\begin{center}
\footnotesize
\begin{tabular}{@{}>{\raggedright\arraybackslash}p{0.22\textwidth}>{\raggedright\arraybackslash}p{0.34\textwidth}>{\raggedright\arraybackslash}p{0.32\textwidth}@{}}
\toprule
半径 & 最小面积余量 & 来源 \\
\midrule
$11323/100000$ & $3797228001573004612011/10^{24}$ & 外部重放证书 \\
$11324/100000$ & $3205113070935301276991/10^{24}$ & \texttt{n100-\allowbreak{}B31-\allowbreak{}at-\allowbreak{}11324.json} \\
$1132301/10^{7}$ & $3791306959813284956819/10^{24}$ & 第 143 阶段余量表 \\
\bottomrule
\end{tabular}
\end{center}

\section{分支 B32}

$B32$ 分为 $M\le25$ 与 $M\ge26$。

\subsection{低 $M$：$81=73+8$ 个整数剖面}

增强的低 $M$ 归约排除 $M\le24$，把 $M25$ 约化为一个唯一的度数为 $5$ 的深点、至多两个高次点，以及四个循环星形模式。随后在全部 $81$ 个整数剖面上评估原始的 $M25$ 费用行：$73$ 行严格为正，$8$ 行为负。这 $8$ 个残差全为 $e0$ 型，带两个度数为 $7$ 的高点、$p\in\{0,2\}$ 以及中心星形类型 $2$ 或 $3$；最差的原始间隙约为 $0.00012804$，这就是为什么只看第一个失败（约 $0.00002125$）会产生误导。

对这个固定的整数域，半侧乘子、面积比与周长罚被重新调为
\[
K=\frac{78619}{283293},\qquad \epsilon=\frac{499143}{911596},\qquad
\tau=\frac{930249}{710647},
\]
其中 FAN 支付严格为 $0$，因为这些行没有 $p3$；并且规范化的边界包络与零斜率铰链在新的 $K$ 下被精确复核（没有任何支付从旧的 $K$ 结转）。精确收据为
\[
\min_{73\ \text{original rows}}=\frac{15600046212127187429}{200000000000000000000000}>0.00007800,
\]
\[
\min_{8\ \text{new rows}}=\frac{1596099436623121073}{12500000000000000000000}>0.00012768 .
\]
一份独立的 \texttt{Fraction} 台账重新合成全部 $81$ 行，并重算 $9$ 个普通面类、$18$ 个外部面类与残余的中心面；所有严格下端点都为正。发现阶段的哨兵与 $K$ 目标错误已被更正，一个把 $K$ 弄混的 \texttt{baseline\_\allowbreak{}numeric\_\allowbreak{}area\_\allowbreak{}margin} 字段通过一份勘误（\texttt{n100-\allowbreak{}discovery-\allowbreak{}erratum-\allowbreak{}145.json}）被正式撤回；接受仅依据精确证书。

\subsection{高 $M$：标记的度数为 $5$ 的点}

高 $M$ 诊断（diagnostic）显示，原始的面/顶点支撑、双正曲率排除以及 $M$/耳/坏面平移检查全部通过，而仅标记低点排除仍有余量
\[
-\frac{490979625631490392613}{10^{24}} .
\]
应对方式是修改穷尽分类，而不是接受这条失败的断言。在保留全部原始平移不等式的前提下，Euler 恒等式与双曲率排除只留下两种可能：
\begin{enumerate}[leftmargin=1.6em,itemsep=2pt]
\item $M26/e0/b0$，全部 $68$ 个内部点的度数为 $6$；
\item $M26/e0/b0$，带一个唯一的度数为 $5$ 的点、一个唯一的度数为 $7$ 的点，以及 $66$ 个度数为 $6$ 的点；度数为 $5$ 的点的标记计数 $p$ 允许取 $0,1,2,3$。
\end{enumerate}
随后按唯一的度数为 $5$ 的点是否拥有一个原始边界邻居来分类：
\begin{itemize}[leftmargin=1.4em,itemsep=3pt]
\item \emph{无原始边界邻居}：此时 $p=0$，原始直接星形支撑适用，$4$ 个高度标记计数乘以 $2$ 个高邻居星形类型给出 $8$ 行，最小余量为
\[
\frac{78693035940218423479}{500000000000000000000000}>0.00015738 .
\]
\item \emph{有原始边界邻居}：此时度数为 $5$ 的点不可能落在任何诱导核心面块的严格内部，因此平面分离与 $/6$ 等周论证仍然适用。核心仍为 $68$ 个顶点、$108$ 个面、一条 $26$ 条边的无弦边界与 $42$ 个深点。深点度超出量 $h$ 为 $0$ 或 $1$，四面度量以最小余量覆盖这两种情形
\[
\frac{225914348193616561751}{500000000000000000000000}>0.00045182 .
\]
\end{itemize}
关键的角台账不依赖于 $p$：设 $\Gamma$ 为标记内部角总和，则环点出现次数为 $(72-h)+52=124-h$，环点圆心角总和为 $(32+h)\pi/3-\Gamma$，原始边界圆心角总和为 $\Gamma-2\pi-26\pi/3$；CAS 断言这两个恒等式以及合并的线性位势恒等式。面计数
\[
\mathrm{DDD}=62+h,\qquad \mathrm{DDR}=20-h,\qquad \mathrm{DRR}=\mathrm{BRR}=26
\]
也只依赖于 $h$。$26$ 对物理 DRR/BRR 半边相互抵消，因此不花费额外预算。

\section{分支 B33}

\subsection{真正的残差：一个唯一的度数为 $5$ 的点}

对 $M\ge26$，若存在度数为 $5$ 的点，则它唯一、类型为 $M26/e0/b0$，且总高曲率为 $2$；因此高点要么是两个度数为 $7$ 的点，要么是一个度数为 $8$ 的点。分类同样依据度数为 $5$ 的点是否有一个原始边界邻居 —— 而报告谨慎地指出 $p=0$ \emph{不}蕴含不存在边界邻居，所以两种情形必须按邻居而非按 $p$ 来划分。

\subsection{无原始边界邻居：一次相干星形替换（coherent star replacement）}

保持有效的 $K=11/25$ 与铰链斜率 $3/100$ 以及全部面费用，唯一的度数为 $5$ 的深点的五个中心面被\emph{一次}替换为一个相干多边形支撑，其中
\[
\epsilon=\frac{26304223}{69960624},\qquad \tau=\frac{110526397}{84434654},
\]
以及一个重新优化的中心位势。在至多两个高点的情形下，邻居二元 $5$-循环只有 $4$ 种计数类型，已由机器检查；原始的 $20$ 个面支撑、$18$ 个外部支撑与 $8$ 个中心支撑被精确重算；没有任何面包含三个不同的高点。每一项循环规范化增益都严格大于 $9/100$，所以旧的负度数为 $5$ 的行加上 $R^2\cdot9/100$ 给出保守的正余量
\[
\frac{185371504379126342167}{250000000000000000000000}\approx0.000741486 .
\]
五个被选中的物理边抵消相反的半边支付；不花费第二个铰链。

\subsection{有原始边界邻居：带标签的环度量}

这里的核心是一个简单的无弦圆盘，$I67/P26/F106$，含 $41$ 个深点。由于 $P26>3$，每个环点至少有一个深邻居、至少有两个核心面；环点还恰好有两个外带 BRR 面与 $p$ 个标记面，因此其原始度数 $d$ 满足 $d-p-2\ge2$，即 $p\le d-4$。这是\emph{这个}核心拓扑的推论，对任意内部点不施加任何约束。因此唯一的环点度数为 $5$ 的点有 $p\in\{0,1\}$；环点度数 $7$ 只需 $p\le3$，环点度数 $8$ 只需 $p\le4$。

在深点超出量 $h\in\{0,1,2\}$ 下，区域面计数由
\[
\mathrm{DDD}+\mathrm{DDR}+26=106,\qquad 3\,\mathrm{DDD}+2\,\mathrm{DDR}+26=246+h,
\]
给出 $\mathrm{DDD}=60+h$ 与 $\mathrm{DDR}=20-h$。$26$ 条核心环边与 $26$ 个外带 BRR 面逐条物理边配对。枚举 $0$、$1$ 或 $2$ 个高点在 $26$-循环上的位置，并聚合五种边类型 RR/RH/HH/RL/HL 连同环点度数 $7$/$8$ 与标记计数，恰好产生 $108$ 个计数模型。选定的有理参数为
\[
LR=-\frac{2217213}{59317159},\quad LH=-\frac{47596282}{80611277},\quad
LL=\frac{20620389}{94559174},\quad LB=-\frac{5596292}{14123693},
\]
\[
\tau=\frac{11510567}{31711524},\qquad K=\frac{17082939}{31059893},
\]
而 $14$ 个全单纯形支撑（DDD、三个 DDR 类、五个 DRR、五个 BRR）在每一项角退化界与定位护栏下被重算。全部 $108$ 行都为正，最小余量
\[
\frac{147611772341741487969}{200000000000000000000000}\approx0.0007380589 .
\]
该分支不花费铰链，也不花费方差。

\subsection{不含度数为 $5$ 的点的组件}

三个全新的子进程重放 $M26$ 七面度量、$M27$ 带弦度量与单耳度量。旧的「$b1/e0$ 不可能」与「$M26$ 中度数为 $7$ 的点不能是深点」两条断言，通过核心连通性/空侧耳生成以及一般的 $/6$ 内部前沿不等式
\[
26+19+\frac{19^2}{6}<106 .
\]
被重新审查。最终独立合成条目重新枚举 $M/e/b/$degree-$5$ 计数，用 CAS 求解 DDD/DDR 恒等式，把全部 $26$-循环签名匹配到 $108$ 个模型，用一份独立的 \texttt{Fraction} 台账从 $14$ 个带符号支撑重新合成 $108$ 个最终面积上界，然后接受。它拒绝在 \texttt{python -\allowbreak{}O} 下运行，其负对照以 $1$ 退出并给出\emph{Assertions must remain enabled}。整个分支报告的总最小正余量仍由低 $M$ 部分决定，约为 $0.00000294$。

早先对统一环位势的一次尝试诚实地失败了：最优 $LR$ 趋于 $0$，前两次有理化都没有留下严格下界；记录的最小预测余量约为 $-0.00067010$，它\emph{未}被接受。只有区分低、高与普通环点的 $14$ 面结构才产生了精确的正证书。

\section{分支 B34}

\subsection{低 $M$：一个孤立的五角星}

两个原始端点曾为负（$M24..25$ 与 $M\ge26$）；诊断阶段只放开 $M24$ 断言以暴露第二个失败，同时让全部 $48$ 个面类、$1455$ 个顶点类、平移不等式、规范化边界、$q4$ 支付与核心前提（precondition）实际通过。低 $M$ 闭合随后使用原始顶点缩减成本与负的 $M$/耳/坏面系数：$M22..23$ 为正；$M25$，或带至少一个耳或坏面的 $M24$，为正；选定一个保证存在的低度数深点后，任何额外的全局度数 $3/4/5$ 点的中性顶点余量覆盖负面积；而若所选深点本身的度数为 $3$ 或 $4$，它自己的顶点余量也覆盖之。剩下的情形是
\[
M24/e0/b0,\quad\text{exactly one global degree-}5\text{ point}
\]
\[
\text{selected as the deep witness},\quad\text{all other interior degrees}\ge6 .
\]
两次剥离随后表明：$I66/M24$ 的诱导核心必是一个简单的 $P24/F106$ 圆盘（$/5$ 排除多个组件、桥与多个块）；度数为 $5$ 的点的邻居全是严格内部核心点；而若它与某个高点相邻，该高点就必须是唯一的度数为 $7$ 的深点，且核心必须是一个无弦简单 $P18/F64$ 圆盘 —— 这与不等式 $18+12+12^2/5<64$ 以及再剥离一次后的 $18+13+13^2/6<64$ 都矛盾。因此度数为 $5$ 的点的全部五个邻居的原始度数为 $6$，五个中心面的位势类型为 $(P,0,0)$。把这五个面\emph{一次}替换为完整的五角星帽
\[
\mathrm{cap}_5=\sqrt{\frac{25\,(25+11\sqrt5)}{32}}
\]
（其中心角和为正 $2\pi$，给出角位势项 $P\pi/3$）产生精确的规范化增益
\[
\frac{70784489700959821434439}{10^{24}}>0.07078448,
\]
以及最终面积余量
\[
\frac{448474306824206627004168866279}{625000000000000000000000000000000}>0.00071755 .
\]
用下界 $1732050807/10^9$ 约束 $\sqrt3$、用上界 $2236067978/10^9$ 约束 $\sqrt5$，并取 $\mathrm{cap}_5<6.23$，一次独立重算复现出正余量
\[
\frac{190482881108285619519923}{292968750000000000000000000}>0.00065018 .
\]

\subsection{高 $M$：三个拓扑槽}

$M26$ 行的正中性低度数顶点余量排除任何度数 $3/4/5$；负的 $M$ 斜率与耳/坏面平移费用排除 $M\ge27$、至少一个耳，以及至少一个坏 $q1$。剩下的情形是
\[
M26/e0/b0,\quad\text{all original interior degrees}\ge6,\quad\text{positive degree excess }2,
\]
其中高点为两个度数为 $7$ 的点或一个度数为 $8$ 的点 —— 并且不假设它们位于核心环上。核心为 $I66/P26/F104$，含 $40$ 个深点；$/6$ 排除多个组件、桥与多个环，留下一个简单圆盘，但 $P26/F104$ 允许一条弦，所以 $B33$ 的无弦结论不能被照搬。设弦数为 $C$、空核心边界三角形数为 $f_3$、$a=C-f_3\ge0$、深点超出量为 $h$，则带弦感知的首次剥离给出
\[
b=20-h-C-3a,\qquad \text{deleted faces}=46-h-C-a,
\]
而 $/6$ 给出
\[
104\ \le\ 46-h-C-a+\frac{(20-h-C-3a)^2}{6}.
\]
在 $0\le h\le2$、$0\le C\le23$、$0\le a\le C$ 上做机器穷举，且前沿非负，恰好留下三个槽：

\begin{center}
\footnotesize
\begin{tabular}{@{}ccccc@{}}
\toprule
核心弦数 $C$ & 深点超出量 $h$ & $a$ & 内部前沿 & 内部面数 \\
\midrule
$0$ & $0$ & $0$ & $20$ & $58$ \\
$0$ & $1$ & $0$ & $19$ & $59$ \\
$1$ & $0$ & $0$ & $19$ & $59$ \\
\bottomrule
\end{tabular}
\end{center}

单弦槽恰好有一个空的核心边界三角形、全部深点的原始度数为 $6$，且环上的总高曲率为 $2$。无弦 $h=1$ 槽有一个度数为 $7$ 的深点与一个度数为 $7$ 的环点；无弦 $h=0$ 槽的全部高曲率都在环上。CAS 发现一般的 DDD 公式漏掉了 $+a$；失败的原脚本被保留，更正后的计数为
\[
\mathrm{DDD}=58+C+h+a,\qquad \mathrm{DDR}=20-C-h-3a .
\]
三个存活槽都有 $a=0$，所以该更正不改变它们。单弦支撑必须保留 $\mathrm{DDD}=59$、$\mathrm{DDR}=19$、$\mathrm{DRR}=25$、配对的 $\mathrm{BRR}=24$、未配对的 $\mathrm{BRR}=2$ 与 $\mathrm{RRR}=1$；把它当作 $26$ 对完全配对的 DRR/BRR 是\emph{不}合法的。

该分支随后以 $23$ 个无弦模型与 $47$ 个单弦模型闭合，二者都被细化到一个由 $10\,000$ 个区间箱组成的全局边界包络。无弦模型使用 $9$ 个带符号面支撑、$26$ 对物理边，铰链/方差支付为零；单弦模型额外区分 $24$ 条配对外部 DRR/BRR 边、一条弦边、一个 RRR 三角形与两个未支付的 BRR 面。独立的 \texttt{Fraction} 重合成与 CAS 圆心角恒等式均通过。该分支的一次统一环位势发现运行记录了约 $-0.00103852$ 的最优预测余量；它只是发现，未接受任何排除。

\section{分支 B35}

诊断保留了全部面、顶点、$q4$ 支付、规范化边界与拓扑前提，并显示原始 $M\le24$ 标记行以及受限的 $M25$ 与 $M26$ 的面积端点为负，而 $M\ge27$、耳、坏 $q1$ 与额外低度数平移余量仍为正。因此严格残差为
\[
M\le24\ \text{flagged low-}M;\qquad M25\ \text{with all interior degrees}\ge6\ (\approx-0.0004727668);
\]
\[
M26\ \text{with all interior degrees}\ge6\ (\approx-0.0002967289).
\]
它返回 \texttt{status = PASS\_\allowbreak{}PREMISES} 与 \texttt{B35\_\allowbreak{}closed = false} —— 正是这个状态词防止把一次前提检查读成一次闭合。

闭合随后完全按 $B34$ 的方式进行：低 $M$ 部分由唯一的未标记度数为 $5$ 的点与核心耳数之间的 $/5$ 矛盾闭合（至多允许一个核心耳）；$M25$ 由 $8$ 个环模型闭合（核心 $P25/F103$、$40$ 个深顶点、$25$ 对物理边、四个度数为 $7$ 的环点且 $pH\in\{11,12\}$）；$M26$ 由 $3$ 个无弦模型加 $5$ 个单弦模型闭合（核心 $P26/F102$、$39$ 个深顶点、$\mathrm{DDD}=56$、$\mathrm{DDR}=[20,0]$）。每个新度量都有自己的完整边界 $K$ 与面求和，并带精确的物理边抵消，且没有任何标量亏空被跨两行支付。合并后的独立证书报告 \texttt{M25\_\allowbreak{}closed = true}、\texttt{M26\_\allowbreak{}closed = true}、\texttt{M\_\allowbreak{}ge27\_\allowbreak{}closed = true}，模型计数为 $\{M25:8,\ M26_{\text{chordless}}:3,\
M26_{\text{onechord}}:5\}$。

\section{尾部分支 $B\ge36$}

\subsection{capped-pair 松弛（capped-pair relaxation）}

对高边界计数，引擎改变了性质。它不再使用全局面台账，而是为每一对相邻的边界胞（boundary cell）构造一个\emph{局部对损失}。在吸收形式中，一个对项具有形状
\[
\Phi(a,b)=\lambda\,\|c_a-c_b\|^2+\frac{H(a)+H(b)}{2}-\frac{\mu\,(\arcsin s_a+\arcsin s_b)}{2},
\]
其中 $a,b$ 是相邻的半弧角，$c_a,c_b$ 是相应的圆心，$H$ 是由弧段与圆心到弧弦的距离构造的几何损失，$\lambda,\mu$ 为罚系数（penalty coefficient）与斜率系数（slope coefficient）。当局部项沿环求和时，使用的是精确恒等式
\[
\sum_i\frac{a_i+a_{i+1}}2=\sum_i a_i=\pi
\]
—— 而不是 Jensen 型的「等弧」猜想。每个局部对的上界随后被送入一个总面积预算，并执行两项标量检查：$B=36$ 处的余量必须严格为正，且每增加一个边界胞的增益必须严格为正，从而把排除扩展到每个 $B\ge36$。

证书脚本实际实现的标量预算是
\begin{equation}\label{eq:scalar-budget}
\underbrace{(1-\mu)\,\pi_{\mathrm{lo}}}_{\text{available arc budget}}
\;-\;\underbrace{\bigl((2m-2-B)\,A+B\,K\bigr)R^2}_{\text{interior and marked triangle area}}
\;-\;\underbrace{B\,L}_{\text{intercept}}\;>\;0,
\qquad
\underbrace{(A-K)R^2-L}_{\text{gain per boundary cell}}\;>\;0,
\end{equation}
其中 $m=100$，$A=129903810568/10^{11}$（$3\sqrt3/4$ 的一个向外有理包围），$K$ 为标记三角形常数。代入冻结常数可精确复现存档的余量，下面两个证书字段正是这样被检查的：
\[
\text{at }R=\tfrac{113237}{10^6},\ L=\tfrac{75097}{10^7}:\quad
\text{margin}=4.906445222464608\times10^{-7},
\]
\[
\text{gain}=5.1549192230158705\times10^{-3},
\]
\[
\text{at }R=\tfrac{22647}{200000},\ L=\tfrac{939}{125000}:\quad
\text{margin}=1.8087101250224213\times10^{-5},
\]
\[
\text{gain}=5.152171860143971\times10^{-3}.
\]
\eqref{eq:scalar-budget} 的底层几何推导位于被吸收的第三方包中；本项目所证明的是该预算的精确算术、下一小节所述有理方形域细分的完备性，以及负对照。

\subsection{有理方形域细分}

对域是一个二维有理方形。它由递归的有理箱细分覆盖；对每个箱，局部对损失被向外界定并与预算比较。对角线上的箱因对称性被跳过并单独计数。三个「$q$-标记侧」三角形支撑被独立验证。被接受的计数为

\begin{center}
\footnotesize
\begin{tabular}{@{}lrrrr@{}}
\toprule
证书 & 半径 & 访问箱数 & 接受箱数 & 对称跳过 \\
\midrule
在 $R_\star=22647/200000$、$L=939/125000$ 处的 $B=36$ & $0.113235$ & $1\,841\,089$ & $917\,297$ & --- \\
在 $R=2831/25000$、$L=19/2500$ 处的 $B\ge37$ & $0.11324$ & $93\,257$ & $46\,038$ & $591$ \\
\bottomrule
\end{tabular}
\end{center}

其中对 $1$、$2$、$3$ 个标记侧的 $q$-支撑箱数分别为 $3761$、$104$ 与 $35$，且对 $B\ge37$ 有 \texttt{pending\_\allowbreak{}boxes} 为空。精确余量为
\[
B=36:\ \frac{1808710125022421182469778967931605940457120574226219}{10^{56}}
\approx1.8087\times10^{-5},
\]
\[
\text{per cell}:\ \frac{2576085930071985489}{5\cdot10^{20}}\approx5.1522\times10^{-3},
\]
\[
B=37:\ \frac{166438291677217339982469778967931605940457120574226219}{10^{56}}
\approx1.6644\times10^{-5},
\]
\[
\text{per cell}:\ \frac{39572580329210081}{7812500000000000000}\approx5.065\times10^{-3}.
\]
$B=36$ 的完整运行耗用 $130.04$ 秒的报告时间；尾部是整个引擎中最昂贵的部分。

\subsection{尾部为何不能延伸到 $0.11324$}

在 $R=2831/25000$ 处，\emph{旧的}固定系数给出的总预算约为 $-0.000160907$：即使其局部支撑仍然有效，尾部也无法在 $0.11324$ 处闭合。固定 $\mu$ 与 $\lambda$ 时，截距（intercept）必须严格低于约 $0.007505530351154887$；试取 $L=0.007505$ 给出约 $+0.0000190926$ 的条件预算，但完整的二维局部支撑从未被验证（一次 $250\,000$ 箱的运行接受了 $124\,491$ 个叶子，留下 $17$ 个待处理栈项，并返回 \texttt{UNDECIDED}）。

这个障碍是精确的，而非数值的。在旧常数下，于有理点
\[
x=y=\frac{9921033}{12500000}
\]
处，对松弛满足
\[
\frac{3754996055077}{500000000000000}\ \le\ \text{pair}\ \le\ \frac{3754996055079}{500000000000000},
\]
因此它的下端点超出旧截距 $1501/200000$ 达
\[
\frac{2496055077}{500000000000000}\approx4.99211\times10^{-6},
\]
并且它也超出 $B=36$ 面积预算所能支付的最大截距
\[
\frac{27019909263463328303182469778967931605940457120574226219}
{3600000000000000000000000000000000000000000000000000000000},
\]
达
\[
\frac{16062333091071696817530221032068394059542879425773781}
{3600000000000000000000000000000000000000000000000000000000}.
\]
因此这个亏空既不能通过提高某个箱预算来修补，也不能通过抬高截距来修补：后者会使 $B=36$ 预算变为负。后来的一次精确诊断表明，在障碍点（obstruction point）处两个支撑项都消失，所以仅仅细化支撑箱在数学上同样不够。障碍证书带有 \texttt{status = PASS\_\allowbreak{}OBSTRUCTION\_\allowbreak{}TO\_\allowbreak{}OLD\_\allowbreak{}RELAXATION}、\texttt{full\_\allowbreak{}B36plus\_\allowbreak{}closed = false}、\texttt{global\_\allowbreak{}bound\_\allowbreak{}improved = false}，并且其 scope 字符串明确声明：这不是一个覆盖构型，也不是对目标下界的反例。

\section{合并器}

最终装配刻意保持单薄，因为一切实质内容都在某个分支证书之内。

\begin{itemize}[leftmargin=1.4em,itemsep=3pt]
\item \texttt{verify/\allowbreak{}out/\allowbreak{}n100-\allowbreak{}candidate-\allowbreak{}complete-\allowbreak{}225.json}（模式 \texttt{n100-\allowbreak{}candidate-\allowbreak{}complete/\allowbreak{}v1}，\texttt{status = PASS}）记录 \texttt{candidate\_\allowbreak{}radius = 22647/\allowbreak{}200000}、\texttt{previous\_\allowbreak{}target\_\allowbreak{}radius = 2831/\allowbreak{}25000}、\texttt{closed\_\allowbreak{}branches = [B31,\dots,B35,B36plus]}、\texttt{all\_\allowbreak{}six\_\allowbreak{}closed = true}、\texttt{global\_\allowbreak{}bound\_\allowbreak{}improved = true}，以散文形式记录单调性论证，并对 $B31$--$B35$ 中的每一个记录源半径 $2831/25000$、继承半径 $22647/200000$ 以及源证书的 SHA-256。尾条目带有它自己的半径与哈希。
\item \texttt{verify/\allowbreak{}out/\allowbreak{}n100-\allowbreak{}six-\allowbreak{}branch-\allowbreak{}ledger-\allowbreak{}226.json}（模式 \texttt{n100-\allowbreak{}six-\allowbreak{}branch-\allowbreak{}ledger/\allowbreak{}v11}）记录 \texttt{closed\_\allowbreak{}branches}、\texttt{pending\_\allowbreak{}branches =
[]}、\texttt{all\_\allowbreak{}six\_\allowbreak{}closed = true}，以及一个 \texttt{scope} 字符串，它用文字说明：只有候选半径被覆盖，且 $R=0.11324$ 仍未被证明，因为它的 $B36$ 直接障碍依然存在。
\item {\footnotesize\texttt{verify/\allowbreak{}out/\allowbreak{}n100-\allowbreak{}small-\allowbreak{}advance-\allowbreak{}controls-\allowbreak{}143.json}} 拒绝两个破坏性情形：一个禁用断言，另一个错配旧半径。
\item {\footnotesize\texttt{verify/\allowbreak{}out/\allowbreak{}n100-\allowbreak{}radius-\allowbreak{}dependency-\allowbreak{}audit-\allowbreak{}143.json}} 追踪 $20$ 个进程与 $28$ 个 \texttt{sys.modules} 注册的源，记录每个验证函数所见的半径，并确认不存在其他全局 $R$；匿名加载的后端由完整的 $626$ 文件哈希清单绑定。
\end{itemize}

\section{分支证书一览}

\begin{center}
\footnotesize
\begin{tabular}{@{}>{\raggedright\arraybackslash}p{0.24\textwidth}>{\raggedright\arraybackslash}p{0.22\textwidth}>{\raggedright\arraybackslash}p{0.22\textwidth}>{\raggedright\arraybackslash}p{0.20\textwidth}@{}}
\toprule
分支 & 完整覆盖范围 & 建立于 & 证书 \\
\midrule
$B31$ & $M\le23$ 与 $M\ge24$，任意耳/坏面 &
$11324/100000$ & \texttt{n100-\allowbreak{}B31-\allowbreak{}at-\allowbreak{}11324.json} \\
$B32$ & $M\le25$（$81=73+8$ 剖面）与 $M\ge26$ &
$2831/25000$ & \texttt{n100-\allowbreak{}B32-\allowbreak{}complete-\allowbreak{}149.json} \\
$B33$ & 低 $M$；$M25$ 相干替换；全部高 $M$ Euler 残差 &
$2831/25000$ & \texttt{n100-\allowbreak{}B33-\allowbreak{}complete-\allowbreak{}167-\allowbreak{}v2.json} \\
$B34$ & $M22..23$、$M24..25$、$M\ge26$（$23$ 无弦 $+$ $47$ 单弦） &
$2831/25000$ & \texttt{n100-\allowbreak{}B34-\allowbreak{}complete-\allowbreak{}180.json} \\
$B35$ & 低 $M$、$M25$、$M26$、$M\ge27$ &
$2831/25000$ & \texttt{n100-\allowbreak{}B35-\allowbreak{}complete-\allowbreak{}203.json} \\
$B\ge36$ & $B=36$ 与 $B=37..99$ &
$22647/200000$（尾部）、$2831/25000$（$B\ge37$） &
\texttt{n100-\allowbreak{}tail-\allowbreak{}candidate-\allowbreak{}224.json}、\texttt{n100-\allowbreak{}tail-\allowbreak{}B37plus-\allowbreak{}210-\allowbreak{}a.json} \\
\bottomrule
\end{tabular}
\end{center}

所有分支共享相同的解析（analytic）依赖，这一点在每个证书中都被陈述，从不隐藏：公共规范化圆心三角剖分；原始边界邻居与严格核心块内部之间的平面分离；一般的 $/6$ 圆盘等周与诱导核心 $/5$ 等周；核心/外带分解；完整星形/周长引理；全单纯形带符号面支撑；以及仅花费一次的物理边抵消（physical edge cancellation）。证书包含整数穷举、CAS 恒等式与精确有理区间算术 —— 它们\emph{不}包含形式化的平面嵌入内核，并且每一个证书都这么说明。

\part*{第五部分 当前前沿：界如何移动，以及它停在哪里}
\addcontentsline{toc}{part}{第五部分 当前前沿：界如何移动，以及它停在哪里}
\stepcounter{partanchor}\setcounter{section}{0}

\section{公认区间及其历史}

当前公认、可机器重放的陈述是
\[
\boxed{\ 0.113235\ <\ \rD(100)\ \le\ 0.1154469548855829\ }
\]
其中下端点由第四部分的六分支（branch）链在 $R_\star=22647/200000$ 处建立，上端点由第二部分的被重放显式构型（configuration）给出。相对间隙 $(U_\star-R_\star)/U_\star$ 约为 $1.92\%$（显示值）。

\subsection{下界轨迹}

\begin{center}
\footnotesize
% 第 1 列必须容纳 $554570693/5000000000$，第 3 列必须容纳最长的
% \texttt{verify/out/...} 路径。宽度须为两个内部间隔留出 \tabcolsep（4\tabcolsep = 24pt）；
% 原先的 0.16/0.24/0.46 溢出 17–21pt。
\begin{tabular}{@{}>{\raggedright\arraybackslash}p{0.20\textwidth}>{\raggedright\arraybackslash}p{0.20\textwidth}>{\raggedright\arraybackslash}p{\dimexpr0.60\textwidth-4\tabcolsep\relax}@{}}
\toprule
下界（lower bound）（精确） & 十进制（显示） & 路线（route） / 证据 \\
\midrule
$1/10$ & $0.100000$ & 仅面积；\texttt{verify/\allowbreak{}out/\allowbreak{}lowerbound\_\allowbreak{}n100.json} (L01--L20) \\
$554570693/5000000000$ & $0.1109141386$ & 修复的整数边路线；\texttt{verify/\allowbreak{}out/\allowbreak{}strengthened-\allowbreak{}lowerbound-\allowbreak{}n100.json} \\
$1111720153/10^{10}$ & $0.1111720153$ & 加细的包络；\texttt{verify/\allowbreak{}out/\allowbreak{}n100-\allowbreak{}algebraic-\allowbreak{}bounds-\allowbreak{}bundle.json} \\
$2787/25000$ & $0.11148$ & 长弦二次亏量；\texttt{verify/\allowbreak{}out/\allowbreak{}n100-\allowbreak{}longchord-\allowbreak{}bundle-\allowbreak{}frozen.json} \\
$11323/100000$ & $0.11323$ & 外部六分支重放；\texttt{verify/\allowbreak{}out/\allowbreak{}external-\allowbreak{}disk100-\allowbreak{}bounds-\allowbreak{}67ebf1c.json} \\
$1132301/10^{7}$ & $0.1132301$ & 同一骨架，半径推进；\texttt{verify/\allowbreak{}out/\allowbreak{}n100-\allowbreak{}new-\allowbreak{}lower-\allowbreak{}bound-\allowbreak{}143.json} \\
$22647/200000$ & $\mathbf{0.113235}$ & 六个分支，尾部在候选半径处重放；\texttt{verify/\allowbreak{}out/\allowbreak{}n100-\allowbreak{}candidate-\allowbreak{}complete-\allowbreak{}225.json} \\
\bottomrule
\end{tabular}
\end{center}

最后两项与前四项性质不同。无构型（configuration-free）的界是单条解析不等式；$0.11323$ 及其后的值来自矛盾引擎（contradiction engine），它们的改进以 $10^{-7}$ 到 $10^{-5}$ 为单位度量，因为每一次推进都要求在新半径处重新闭合每一个分支。

\subsection{上界轨迹}

\begin{center}
\footnotesize
% \mathbf{0.1154469548855829} 是第 2 列最宽的一项，故其框宽于第 1 列；
% 第 3 列取扣除 4\tabcolsep 之后的余量。
\begin{tabular}{@{}>{\raggedright\arraybackslash}p{0.20\textwidth}>{\raggedright\arraybackslash}p{0.26\textwidth}>{\raggedright\arraybackslash}p{\dimexpr0.54\textwidth-4\tabcolsep\relax}@{}}
\toprule
上界（upper bound） & 十进制（显示） & 路线 / 证据 \\
\midrule
$153/1250$ & $0.1224$ & 历史构型的 stage-1 严格证书（certificate）；\texttt{verify/\allowbreak{}out/\allowbreak{}coverage\_\allowbreak{}n100.json} \\
$1223433/10^{7}$ & $0.1223433$ & stage-2 服务器搜索，六个候选；\texttt{verify/\allowbreak{}out/\allowbreak{}pilot\_\allowbreak{}n100\_\allowbreak{}stage2\_\allowbreak{}summary.json} \\
$1/\sqrt{67}$ & $0.1221694443563052$ & 三角格 $97{+}3$；\texttt{verify/\allowbreak{}out/\allowbreak{}lattice97\_\allowbreak{}upperbound\_\allowbreak{}n100.json} \\
$18761/156250$ & $0.1200704$ & 协作第 1 轮；\texttt{n100-\allowbreak{}cooperative-\allowbreak{}anchor-\allowbreak{}tight/\allowbreak{}replay.json} \\
$119891/10^{6}$ & $0.119891$ & 精化；\texttt{n100-\allowbreak{}cooperative-\allowbreak{}refine-\allowbreak{}tight/\allowbreak{}replay.json} \\
$119877/10^{6}$ & $0.119877$ & 微调；\texttt{n100-\allowbreak{}cooperative-\allowbreak{}polish-\allowbreak{}tight/\allowbreak{}replay.json} \\
$1154469548855829/10^{16}$ & $\mathbf{0.1154469548855829}$ & 吸收的外部构型；\texttt{verify/\allowbreak{}out/\allowbreak{}external-\allowbreak{}disk100-\allowbreak{}margin-\allowbreak{}67ebf1c.\allowbreak{}json} \\
\bottomrule
\end{tabular}
\end{center}

余量证书值得仔细阅读，因为它表明这一改进是精确的而非表面上的。被冻结的 $100$ 中心有理构型经检验覆盖整个闭单位圆盘，其
\[
\text{covering squared radius}\ \le\
\tfrac{1332799939235381342323027277241}{10^{33}}
=0.013327999392353813\ldots=U_\star^2 ,
\]
由检查 $128\,436$ 个内部三元组、$9\,900$ 对边界对、$100$ 对对径点对与 $1$ 个例外情形（共 $138\,437$ 个对象，精确整数算术）得到。证书记录了一个\emph{参考}界 $115446954885583/10^{15}=0.115446954885583$ 与本项目的局部界 $1154469548855829/10^{16}$；两者之差恰好为 $1/10^{16}$。所以公认的上端点是对参考值的 $10^{-16}$ 的严格改进，而不是对它的重述。

\section{从 $0.11323$ 到 $0.1132301$：移植与第一次推进}

Stage 139 把计划从接口（interface）整合切回半径推进。它在一个隔离的工作副本内重放了吸收来的六分支骨架（退出码 $0$，绑定 $626$ 个文件，运行后源与副本的哈希不变），并产生了一个六端点有理半径扫描（$0.11323$，$0.113231$，$0.113235$，$0.11324$，$0.11325$，$0.1133$），对每个端点记录计数桥（counting bridge）、$B31$ 固定费用行与尾部的条件总预算。浮现出两个事实：

\begin{itemize}[leftmargin=1.4em,itemsep=2pt]
\item 在 $R=0.11324$ 处，$99$ 圆盘排除仍然成立，$100$ 个圆盘的必要边界计数仍为 $31..39$，并且 $B31$ 末行、耳（ear）/坏面（bad face）符号、度数预算与 fan 必要限制全部通过。在 $R=0.11325$ 处，$B31$ 的 fan 前提（precondition）不再闭合，因此整个分支在那里失败，尽管其最终面积余量仍为正。
\item 在旧尾部常数下，总预算只允许半径达到某个阈值，该阈值位于
\[
\Bigl[\frac{11323679462887}{10^{14}},\ \frac{113236794628871}{10^{15}}\Bigr),
\]
之中；并且在 $R=0.11324$ 处旧尾部预算约为 $-0.000160907$。在固定 $\mu$ 与 $\lambda$ 的情况下，截距（intercept）必须降到约 $0.007505530351154887$ 以下。
\end{itemize}

随后在 $R=11324/100000$ 处显式重放了 $B31$ 的算术
（\texttt{n100-\allowbreak{}B31-\allowbreak{}at-\allowbreak{}11324.json}，退出码 $0$，$40$ 个面支持与 $32$ 个顶点支持，最小面积余量 $3205113070935301276991/10^{24}\approx0.003205113$），其中固定的无量纲费用/势/支持系数保持不变，只调整了半径（以及提案中的半径元数据）。证书把几何转移义务与算术重放分开列出，这是记录「算术仍然闭合」不等于「几何仍然适用」的诚实方式。

Stage 143 随后把半径推进到 $1132301/10^{7}$，并在原始系数未改动的情况下重新闭合了全部六个分支。同半径的精确余量为

\begin{center}
\footnotesize
\begin{tabular}{@{}ll@{}}
\toprule
分支 & 在 $R=1132301/10^{7}$ 处的最小面积余量 \\
\midrule
$B31$ & $3791306959813284956819/10^{24}$ \\
$B32$ & $40052428178643350631/125000000000000000000000$ \\
$B33$ & $141216899554184842489/10^{24}$ \\
$B34$ & $57252253086154121097/250000000000000000000000$ \\
$B35$ & $37566017111533099749/500000000000000000000000$ \\
$B\ge36$ & $33605092614133890926309778967931605940457120574226219/10^{56}$ \\
\bottomrule
\end{tabular}
\end{center}

其中尾部访问 $3\,397\,109$ 个盒并接受 $1\,694\,095$ 个，且每增加一个 $B$ 的余量严格为正。新的归一化中心范数平方上界为 $98717894445399/10^{14}$；同一半径处的计数桥排除了所有 $m\le99$，对 $m=100$ 只留下 $B31$--$39$。改进恰好是 $1/10^{7}$，并且该阶段明确说明它既没有使用 stage-141 的 $M24$ 强化，也没有对铰链斜率（slope）、费用系数或支持行作任何改动。

\section{在 $0.11324$ 处逐个闭合分支}

在 stage 144 到 211 之间，六个分支在更困难的目标 $R=2831/25000=0.11324$ 处被逐个重新闭合，失败也保留在记录中。

\begin{center}
\footnotesize
\begin{tabular}{@{}>{\raggedright\arraybackslash}p{0.16\textwidth}>{\raggedright\arraybackslash}p{0.24\textwidth}>{\raggedright\arraybackslash}p{0.46\textwidth}@{}}
\toprule
阶段 & 分支 & 结果 \\
\midrule
144--149 & $B32$ & 完成：$M\le25$（$81=73+8$ 个 profile）与 $M\ge26$；发现并修复了两处真实失败（发现层（discovery layer）字段中一个被弄混的 $K$，已以勘误撤回；以及一个被标记为低的排除，它必须被新的分类取代而不能被接受） \\
150--168 & $B33$ & 完成：低 $M$、相干的 deep-$5$ 星以及 $108$ 个带标号环模型；第一个均匀环势以约 $-0.00067010$ 的预测余量失败，未被接受 \\
169--185 & $B34$ & 完成：低 $M$（孤立五星）与高 $M$（$23$ 个无弦 $+$ $47$ 个单弦模型）；两个负端点必须先由诊断（diagnostic）暴露出来才能修复 \\
186--211 & $B35$ & 完成：低 $M$、$M25$（$8$ 个环模型）、$M26$（$3$ 个无弦 $+$ $5$ 个单弦）、$M\ge27$ \\
209--211 & $B\ge37$ & 在 $R=2831/25000$、截距 $19/2500$ 下完成：访问 $93\,257$，接受 $46\,038$，对称跳过 $591$，待处理 $0$ \\
206--229 & $B=36$ & \textbf{未闭合}：精确障碍，随后是一个在更小半径处确实闭合的保守候选 \\
\bottomrule
\end{tabular}
\end{center}

正是在这里，公认界成为 $0.113235$ 而不是 $0.11324$：六个分支中有五个在更困难的半径处可用，但尾部不可用，而一条链的强度只等于它最弱的分支。

\section{$B=36$ 的障碍，以及为什么显而易见的修复都失败}

\subsection{精确障碍}

在 $R=2831/25000$ 与冻结常数（$\lambda=3/50$，$\mu=92531/10^{7}$，标记常数 $K=973/3125$，旧截距 $L=1501/200000$）下，有理点
\[
x=y=\frac{9921033}{12500000}
\]
的一个 pair 松弛的区间为
$[3754996055077/5\cdot10^{14},\,3754996055079/5\cdot10^{14}]$。其下端点超出旧截距 $2496055077/5\cdot10^{14}\approx4.99211\times10^{-6}$，并超出 $B=36$ 面积允许的最大截距
$16062333091071696817530221032068394059542879425773781/(36\cdot10^{55})$。因此障碍证书返回 \texttt{PASS\_\allowbreak{}OBSTRUCTION\_\allowbreak{}TO\_\allowbreak{}OLD\_\allowbreak{}RELAXATION}，伴随
\texttt{full\_\allowbreak{}B36plus\_\allowbreak{}closed = false} 与 \texttt{global\_\allowbreak{}bound\_\allowbreak{}improved = false}，其 scope 字符串声明：这不是一个覆盖构型，也不是目标下界的反例。

\subsection{四种无效的修复}

\begin{enumerate}[leftmargin=1.6em,itemsep=3pt]
\item \textbf{更多盒。} 障碍是一个精确有理点上的逐点不等式；该点处的 pair 值就是它本身，所以细分无法把它降低。后来的一次精确诊断确认两个支持项在障碍点处都消失
（\texttt{n100-\allowbreak{}tail-\allowbreak{}point-\allowbreak{}branch-\allowbreak{}212.json}，
\texttt{n100-\allowbreak{}tail-\allowbreak{}direct-\allowbreak{}components-\allowbreak{}216.json}），因此细化支持盒在数学上是不够的。
\item \textbf{更高的截距。} 提高 $L$ 会使 $B=36$ 预算为负，因为该预算在旧截距处已经失败。
\item \textbf{复用 $B\ge37$ 的截距。} $B\ge37$ 证书在 $L=19/2500$ 处闭合；把它用在 $B=36$ 上会使该预算为负。
\item \textbf{斜率与罚系数（penalty coefficient）的同步重调。} 在直接 pair/角族上做更广的发现层搜索得到了最优值
\[
\lambda=0.05944697773645484,\quad \mu=0.010360567875057177,\quad
K\approx0.31268174443157726,
\]
其局部峰值约为 $0.00739518789840696$，而 $B=36$ 允许的截距约为 $0.007391936425058569$：每胞亏量约为 $-3.25147334839\times10^{-6}$，$B=36$ 总亏量约为 $-1.17053040542\times10^{-4}$。这与同一族中的一次独立搜索一致。因此在冻结的直接 pair/角松弛内，简单的同步重调无法闭合 $B=36$。
\end{enumerate}

\subsection{实际尝试过并中止的工作}

两个有理候选被启动但没有跑完：

\begin{itemize}[leftmargin=1.4em,itemsep=2pt]
\item $R=113237/10^{6}$ 配 $L=75097/10^{7}$：标量预算为正（余量 $4.906445222464608\times10^{-7}$，增益 $5.1549192230158705\times10^{-3}$），并且在 $1132301/10^{7}$ 与 $2831/25000$ 之间的 $101$ 点有理扫描表明，已知的障碍点（obstruction point）直到精确样本 $R=11323703/100000000=0.11323703$ 仍在预算内，而从精确样本 $R=113237129/1000000000$ 起严格超出（$101$ 个样本中有 $30$ 个被该单点阻塞）。该扫描自身的 scope 字段写明它只是诊断：未被该点阻塞的样本并不是全局 $B36$ 证书。预算边缘附近的全 pair 域重放细分过久，被中止。
\item $R=113236/10^{6}$ 配 $L=0.007511$，以及更保守的 $L=0.00751$：两者因同样的原因被中止。
\end{itemize}

两者都没有被报告为失败，也都不是证书。保守候选 $R=22647/200000=0.113235$ 配 $L=7512/1000000=939/125000$ 拥有宽得多的预算余量，并且确实跑完了：访问 $1\,841\,089$ 个盒，接受 $917\,297$ 个，三个 $q$-支持以 $3761/104/35$ 个盒通过，$B=36$ 余量与每 $B$ 增益都严格为正。这就是被接受的候选，也是前沿停在 $0.113235$ 的原因。

\begin{remark}[精确陈述的未决窗口]
可验证的未决窗口是 $(0.113235,\,0.11324]$。在它内部，固定旧参数的单点扫描表明，从 $113237129/1000000000$ 起，仅那个特定的障碍点就阻塞了旧松弛——但这是一个关于单点处单个松弛的陈述，\emph{不是}对所有新路线的数学上界。
\end{remark}

\section{路线目录所尊重的三条硬边界}

第六部分的系统目录明确了三条边界，值得在目录本身之前先说明：

\begin{enumerate}[leftmargin=1.6em,itemsep=3pt]
\item 递归的内圆盘计数（先 $\ge56$，然后 $\ge58$ 个内圆盘，于是 $k\le44$，再 $k\le42$ 个服务边界的圆盘）是计数论证的推论，\emph{不是}新的全局半径界。在所有这些阶段中，全局区间始终未变。
\item 进展 $0.11516023\to0.11521651$ 是条件分支 $k\ge55$ \emph{内部}的阈值推进。互补分支 $k\le54$ 从未沿这条线被排除，某个阶段甚至给出了聚合松弛的一个显式整数可行点（$34$ 个近端 $+$ $20$ 个尾端 $+$ $46$ 个非服务 $=100$ 个圆盘，其中 $54$ 个服务），这禁止该松弛的任何逻辑推论为空。
\item 接口/契约阶段中的每一个 \texttt{PASS} 都意味着契约登记、重放一致性或计时采样——绝不是几何闭合。全部 $174$ 个 legacy 源都是 \texttt{PRECONDITION\_\allowbreak{}MISSING}，依赖门是 \texttt{UNDECIDED}。
\end{enumerate}

\section{一次真正的推进需要什么}

前沿不是被算力阻塞的，而是被一条不等式的强度阻塞的。证书与诊断指向同一个结论：

\begin{itemize}[leftmargin=1.4em,itemsep=3pt]
\item \textbf{$B=36$ 的一条新几何不等式}：一个共同端点透镜（lens）约束，或一个真正耦合的多边平均不等式，用以取代直接 pair/角松弛。仅仅重新表达 pair 表达式中已有某一项的改写不会增加信息——这一点已被显式测试（把粗糙的 $4\lambda R^2$ 距离帽替换为精确中心距不会带来任何新东西，因为该项已经在原表达式中）。
\item \textbf{一个可续跑的 pair 盒队列（resumable queue）}：逐轮持久化待处理盒，以哈希绑定参数，缓存（caching）反三角与局部支持区间，并用多进程把域分片（sharding）为对角与非对角区域。在 $0.113236$ 与 $0.113237$ 处被中止的运行失败于墙钟时间，而不是数学；没有队列，每个新候选都要从根重新开始。
\item \textbf{一个更锐的普适界}（若要绕过引擎）：无构型最佳值是 $0.11148$，比公认的引擎界低 $1.55\%$，因此它不是近期的替代品。
\item \textbf{精确值路线}：完整枚举内部映射族，处理 $m<100$、退化构型，并把两个界闭合到同一个代数数上，随后做不可约性与根隔离。从第三方材料吸收的平方松弛 / 实量词消去理论是唯一已知能终止并给出精确值的路线，但它没有可用的 $n=100$ 后端：被编译的多项式系统的零点集处处奇异，所以正则 Krawczyk 盒无法替代完整的 QE 后端，而变量预算（$605$ 个几何 $+$ $1318$ 个松弛 $=1923$，其中 $1922$ 个被消去）是预算，不是运行时间。
\end{itemize}

\part*{第六部分 路线目录：每个族、其构造与状态}
\addcontentsline{toc}{part}{第六部分 路线目录：每个族、其构造与状态}
\stepcounter{partanchor}\setcounter{section}{0}

\noindent
本部分是全文其余内容所准备的系统清册。每一行给出路线（route）、其构造、其状态与其证据。状态词表：\textbf{闭合}（closed，在其声明的域内存在严格证书（certificate）——可能是条件子分支）、\textbf{部分排除}（partial，真实的严格排除只覆盖声明域的一部分）、\textbf{仅诊断}（diagnostic，只有测量、审计或瓶颈）、\textbf{证伪}（falsified，该断言或其前提被显式驳倒，或存在反例）、\textbf{未决}（undecided，跑到完成但零新增排除）、\textbf{接口登记}（interface，\texttt{PASS} 意味着契约登记或重放一致性，\emph{不是}几何闭合）、\textbf{上界构型}（configuration，一个上界（upper bound）构造）。除另有说明外，本部分的每条路线都带有 \texttt{machine\_\allowbreak{}verified\_\allowbreak{}geometry = false} 与 \texttt{global\_\allowbreak{}optimality = false}，并且那个时代的报告所引用的全局区间是 $11323/100000<\rD(100)\le1154469548855829/10^{16}$。制品路径相对于仓库根给出；它们中哪些位于已发布镜像中，见第七部分。

\section{组合与拓扑路线}

这些路线的目的是把连续覆盖变成有限组合对象：平面映射、dart map、带固定面目录的三角剖分，或一个边界字。

\begin{itemize}[leftmargin=1.2em,itemsep=3pt]
\item \textbf{平面 / dart-map 编译器。} 具有独立边标识的顶点、面置换、有序四轴标记；半边编码 $(\alpha,\phi)$ 保持多重边标识；界 $V\le2m+2$，$E\le3m+1$；模型预算 $k_M=19m+22$（$n=100$ 时为 $1922$）\emph{状态：仅诊断（diagnostic）。} 证据：\texttt{verify/\allowbreak{}n100\_\allowbreak{}generic\_\allowbreak{}map\_\allowbreak{}compiler.py}, \texttt{n100\_\allowbreak{}dart\_\allowbreak{}map\_\allowbreak{}canonicalization.py}, \texttt{n100\_\allowbreak{}raw\_\allowbreak{}dart\_\allowbreak{}model\_\allowbreak{}adapter.py} \\
映射规范化
\end{itemize}


\section{边界圆、轴与盒剪枝（box pruning）路线}

这一族直接攻击构型空间：传播中心坐标上的区间约束、分裂盒，并计数有多少圆盘可以服务边界。

\begin{itemize}[leftmargin=1.2em,itemsep=3pt]
\item \textbf{容量（capacity）前沿。} 边界服务上的容量 / Hall 型必要条件；盒感知剪枝；非分离三点容量切割 \emph{状态：部分排除（partial）。} 证据：{\footnotesize\texttt{docs/\allowbreak{}result-\allowbreak{}n100-\allowbreak{}stage25-\allowbreak{}capacity-\allowbreak{}frontier.md}} \\
盒感知的共中心剪枝
\end{itemize}


\section{弧对、禁配（nogood）与学习路线}

这一族在一个固定核心（core）与一个收窄的父域上工作，学习边界瓦片标号之间的必要关系，并追问还剩多少支持。

\begin{itemize}[leftmargin=1.2em,itemsep=3pt]
\item \textbf{弧对一致性与环支持。} 共享一个中心的弧瓦片之间的必要 pair 关系（顶点盒最小距离 $\le2U$）；$289\,530/277\,368$ 对约化为 $1866/2312$ 个标号；通过 $32$ 顶点环上的独立集 DP 得到继承的有限成对模型的精确支持 \emph{状态：未决（undecided）。} 证据：\texttt{verify/\allowbreak{}out/\allowbreak{}n100-\allowbreak{}arc-\allowbreak{}pair-\allowbreak{}consistency.json}, \texttt{n100-\allowbreak{}arc-\allowbreak{}cycle-\allowbreak{}support.json} \\
有序与单调环
\end{itemize}


\section{透镜（lens）、对偶、递归计数与接口路线}

\begin{itemize}[leftmargin=1.2em,itemsep=3pt]
\item \textbf{加细的透镜上界。} 在 $d_0=1-r$ 处对交透镜的统一包围，配合 $\min(2rb-b^2,\ s^2-(1-2r+a)^2)$ 的逐条带上 Darboux 和：$256$ 条带给出 $0.11521249$，$4096$ 条带给出 $11521651/10^{8}$（$\delta=101/10000$），把已认证的 $k\ge55$ 半径从 $0.11516023$ 提升到 $0.11521651$ \emph{状态：闭合（closed，条件分支）。} 证据：{\footnotesize\texttt{docs/\allowbreak{}result-\allowbreak{}n100-\allowbreak{}stage94-\allowbreak{}refined-\allowbreak{}lens.md}} \\
联合支持（已被取代）
\end{itemize}


\section{能量、力与局部排除路线}

\begin{itemize}[leftmargin=1.2em,itemsep=3pt]
\item \textbf{固定应力的调和角能量。} $E(q)=\sum c_{ij}\|q_i-q_j\|^2$，其中 $v^{\mathsf T}H_{\text{ang}}v=2\sum c_{ij}(q_i\cdot q_j)(v_i-v_j)^2$；曾希望一个 Loewner 界能在角域上给出凸性 \emph{状态：证伪（falsified）。} 证据：{\footnotesize\texttt{verify/\allowbreak{}out/\allowbreak{}n100-\allowbreak{}angular-\allowbreak{}energy-\allowbreak{}independent.json}}, \texttt{n100-\allowbreak{}angular-\allowbreak{}energy-\allowbreak{}negative-\allowbreak{}witness.json} \\
角能量，两个精确障碍
\end{itemize}


\section{半径扫描与工程化路线}

\begin{itemize}[leftmargin=1.2em,itemsep=3pt]
\item \textbf{六端点半径扫描。} 在 $0.11323$，$0.113231$，$0.113235$，$0.11324$，$0.11325$，$0.1133$ 处的精确有理扫描，记录计数桥、$B31$ 固定费用行与尾部的条件总预算；在 $0.11325$ 处 $B31$ 的 fan 前提已经失败 \emph{状态：仅诊断（diagnostic）。} 证据：{\footnotesize\texttt{verify/\allowbreak{}out/\allowbreak{}n100-\allowbreak{}lower-\allowbreak{}radius-\allowbreak{}bottleneck-\allowbreak{}firstpass.json}} \\
$101$ 点有理扫描
\end{itemize}

\part*{第七部分：复现、产物与依赖}
\addcontentsline{toc}{part}{第七部分：复现、产物与依赖}
\stepcounter{partanchor}\setcounter{section}{0}

\section{入口点}

\subsection{发表一致性检查}

最便宜的单一入口点是
\begin{verbatim}
python reproduce.py          # from the repository root
\end{verbatim}
它加载四份核心证书（certificate），要求它们的状态，要求
\texttt{candidate\_\allowbreak{}radius = 22647/\allowbreak{}200000}、\texttt{accepted\_\allowbreak{}global\_\allowbreak{}lower\_\allowbreak{}bound = 22647/\allowbreak{}200000}、
\texttt{all\_\allowbreak{}six\_\allowbreak{}closed = true} 以及 \texttt{pending\_\allowbreak{}branches = []}，要求字符串
「0.11324」仍然出现在 README 的证据边界一节中，并为每份证书打印 SHA-256。它\emph{不}重新运行历史重算，并且它明确说明这一点。

它在本次审计（audit）中的输出是
\begin{verbatim}
n=100 publication consistency: PASS
accepted bound: r_D(100) > 0.113235
0.11324: not closed; B36 residual retained
\end{verbatim}
其后每份证书一行。证书名称及其 SHA-256 摘要（digest）按打印原样（此处仅为页面宽度而折行）为：
\begin{itemize}[leftmargin=1.4em,itemsep=2pt]
\item {\footnotesize\texttt{verify/\allowbreak{}out/\allowbreak{}n100-\allowbreak{}candidate-\allowbreak{}complete-\allowbreak{}225.json}}\\
{\footnotesize\texttt{c58c71562ed014c0f716f02137dcf473d98fd73e6888f89aa9d06f08a41707a5}}
\item {\footnotesize\texttt{verify/\allowbreak{}out/\allowbreak{}n100-\allowbreak{}tail-\allowbreak{}candidate-\allowbreak{}224.json}}\\
{\footnotesize\texttt{b5ca9c60be73136356c5c77a3b13a9d9c22a07d265299cb3acda2262a5715fee}}
\item {\footnotesize\texttt{verify/\allowbreak{}out/\allowbreak{}n100-\allowbreak{}candidate-\allowbreak{}budget-\allowbreak{}222.json}}\\
{\footnotesize\texttt{41c0106a015c31c7a323639368cd9ff241c5d745c8a4f842356647dd4922d0e2}}
\item {\footnotesize\texttt{verify/\allowbreak{}out/\allowbreak{}independent\_\allowbreak{}lattice97\_\allowbreak{}upperbound\_\allowbreak{}n100.json}}\\
{\footnotesize\texttt{c9f07ee0be9705cea9349fba2023af3a70824ee791373910809c802b2b1b4f4d}}
\end{itemize}

\subsection{重新运行各条路线}

每条路线都有自己的入口点，而每个入口点都拒绝覆盖已存在的输出目标，因此真正的重跑需要一条新路径。报告所记录的代表性命令如下：

\begin{center}
\footnotesize
\begin{tabular}{@{}>{\raggedright\arraybackslash}p{0.42\textwidth}>{\raggedright\arraybackslash}p{0.44\textwidth}@{}}
\toprule
命令 & 路线 \\
\midrule
\texttt{python verify/\allowbreak{}lattice97\_\allowbreak{}upperbound\_\allowbreak{}n100.py} & 三角晶格上界 \\
{\footnotesize\texttt{python verify/\allowbreak{}independent\_\allowbreak{}lattice97\_\allowbreak{}upperbound\_\allowbreak{}n100.py}} & 独立上界重放 \\
\texttt{python verify/\allowbreak{}run\_\allowbreak{}algebraic\_\allowbreak{}bounds\_\allowbreak{}n100.py --out <fresh>} & 加强包络 + 晶格候选 \\
{\footnotesize\texttt{python verify/\allowbreak{}run\_\allowbreak{}strengthened\_\allowbreak{}lowerbound\_\allowbreak{}n100.py --out <fresh>}} & 整数边计数界 \\
\texttt{python verify/\allowbreak{}run\_\allowbreak{}longchord\_\allowbreak{}n100.py --out <fresh>} & 长弦界 \\
\texttt{python verify/\allowbreak{}run\_\allowbreak{}n100\_\allowbreak{}delivery\_\allowbreak{}audit.py --run-\allowbreak{}dir <fresh>} & 封存与快速重放 \\
\texttt{python verify/\allowbreak{}run\_\allowbreak{}n100\_\allowbreak{}interface\_\allowbreak{}suite.py --run-\allowbreak{}dir <fresh>} & 九批接口重放 \\
{\footnotesize\texttt{python verify/\allowbreak{}run\_\allowbreak{}n100\_\allowbreak{}boundary\_\allowbreak{}line\_\allowbreak{}replay.py --run-\allowbreak{}dir <fresh>}} & $m99$ 边界线 \\
\texttt{python verify/\allowbreak{}run\_\allowbreak{}n100\_\allowbreak{}full\_\allowbreak{}revalidation.py --run-\allowbreak{}dir <fresh>} & 全部 $4776$ 个构型 \\
\texttt{python verify/\allowbreak{}run\_\allowbreak{}n100\_\allowbreak{}interface\_\allowbreak{}timing.py --run-\allowbreak{}dir <fresh>} & 计时校准 \\
{\footnotesize\texttt{python verify/\allowbreak{}run\_\allowbreak{}n100\_\allowbreak{}composition\_\allowbreak{}scaling.py --run-\allowbreak{}dir <fresh>}} & 组合与缩放 \\
{\footnotesize\texttt{python verify/\allowbreak{}run\_\allowbreak{}n100\_\allowbreak{}shared\_\allowbreak{}support\_\allowbreak{}timing.py --run-\allowbreak{}dir <fresh>}} & 共享支撑会话 \\
\bottomrule
\end{tabular}
\end{center}

其中三条（计时、组合/缩放与共享支撑）写入不可覆盖的最终目标，因而\emph{未}被统一地做输出隔离；整合报告把这记录为一项尚欠的包装器债务，而不是把它藏起来。分支（branch）层面的重放由
\texttt{verify/\allowbreak{}n100\_\allowbreak{}B31\_\allowbreak{}new\_\allowbreak{}radius.py}、\texttt{n100\_\allowbreak{}B32\_\allowbreak{}complete\_\allowbreak{}149.py}、
\texttt{n100\_\allowbreak{}B33\_\allowbreak{}complete\_\allowbreak{}167\_\allowbreak{}v2.py}、\texttt{n100\_\allowbreak{}B34\_\allowbreak{}complete\_\allowbreak{}180.py}、
\texttt{n100\_\allowbreak{}B35\_\allowbreak{}complete\_\allowbreak{}203.py}、\texttt{n100\_\allowbreak{}tail\_\allowbreak{}B37plus\_\allowbreak{}resumable\_\allowbreak{}210.py}、
\texttt{n100\_\allowbreak{}tail\_\allowbreak{}candidate\_\allowbreak{}224.py}、\texttt{n100\_\allowbreak{}candidate\_\allowbreak{}complete\_\allowbreak{}225.py} 以及
\texttt{n100\_\allowbreak{}candidate\_\allowbreak{}ledger\_\allowbreak{}226.py} 驱动。$B33$ 与 $B34$ 的入口点要求断言处于启用状态：它们拒绝 \texttt{python -\allowbreak{}O} 并报告\emph{断言必须保持启用}。

\section{产物覆盖：已发布镜像中实际有什么}

本文档是针对工作仓库撰写的。由于一项可复现性声明的价值只取决于它所指向的字节，我们为本专著运行了一次显式的产物（artifact）覆盖审计，其结果属于论文的一部分。

\begin{center}
\footnotesize
\begin{tabular}{@{}>{\raggedright\arraybackslash}p{0.28\textwidth}>{\raggedright\arraybackslash}p{0.60\textwidth}@{}}
\toprule
发现 & 细节 \\
\midrule
镜像（mirror）是一个子集 & 在镜像自身的 README、状态文件与阶段报告所引用的 $100$ 条产物路径中，只有 $19$ 条存在于已发布镜像 \texttt{n100-\allowbreak{}public/\allowbreak{}} 之下；$81$ 条在那里缺失，而这 $81$ 条中有 $80$ 条确实存在于父工作仓库中（有一条探测路径是凭空编造的名字，不计入该计数） \\
\texttt{proofs/\allowbreak{}} 缺失 & \texttt{n100-\allowbreak{}public/\allowbreak{}proofs/\allowbreak{}} 不存在，尽管两份晶格证书都引用了
\texttt{proofs/\allowbreak{}n100-\allowbreak{}lattice97-\allowbreak{}algebraic-\allowbreak{}upperbound.\allowbreak{}md}。
该文件存在于父仓库中，其 SHA-256
\texttt{73d39456b70fdcae\allowbreak{}41ede5718a608d13\allowbreak{}9661a6c7f2a6\allowbreak{}26cf95b4e348be19\allowbreak{}600b}
与主晶格证书所记录的 \texttt{proof\_\allowbreak{}sha256} 字段相符 \\
移植包 & \texttt{runtime/\allowbreak{}disk100-\allowbreak{}release-\allowbreak{}lf-\allowbreak{}67ebf1c/\allowbreak{}} 不在镜像中，但存在于父仓库中，其中包括分支证书所引用的全部四个源码
(\texttt{verify\_\allowbreak{}B31\_\allowbreak{}global\_\allowbreak{}low\_\allowbreak{}degree\_\allowbreak{}0p11323.py},
\texttt{verify\_\allowbreak{}capped\_\allowbreak{}pair\_\allowbreak{}B36plus\_\allowbreak{}0p11323.py},
{\footnotesize\texttt{pair\_\allowbreak{}bound\_\allowbreak{}certificate\_\allowbreak{}20261004/\allowbreak{}verify\_\allowbreak{}pair\_\allowbreak{}bound.py}},
\texttt{threshold\_\allowbreak{}11305\_\allowbreak{}components/\allowbreak{}verify\_\allowbreak{}remaining\_\allowbreak{}cases.py})。因此仅从镜像出发的重放无法执行分支检查器 \\
中间证书 & 镜像保留了最终的组合证书与尾部证书，但许多中间子证书（例如 $B32$/$B33$/$B34$/$B35$ 的分量回执与 stage-143 的装配回执）缺失；它们的数字此后只能从阶段报告中引用，而不能从 JSON 中重新读取 \\
仅根目录产物 & \texttt{verify/\allowbreak{}out/\allowbreak{}external-\allowbreak{}disk100-\allowbreak{}local-\allowbreak{}quadtree/\allowbreak{}replay.json} 与
{\footnotesize\texttt{verify/\allowbreak{}out/\allowbreak{}external-\allowbreak{}disk100-\allowbreak{}strict-\allowbreak{}local-\allowbreak{}67ebf1c.json}} 被镜像的状态文件引用，但只存在于父仓库中 \\
\bottomrule
\end{tabular}
\end{center}

\begin{remark}[对本文档的后果]
第二至第六部分中的每一条定量断言都是对照存在于父仓库中的产物核对的，而不是对照镜像核对的。若某个数字只能取自阶段报告，因为其底层证书不在镜像中，正文会说明这一点或引用该报告。这正是该计划对其自身各阶段所施加的同一条纪律，被施加到发表层上。
\end{remark}

\section{本次审计发现的绑定不一致}

发现了三处绑定问题，它们被记录在此，而不是被悄悄修复。

\begin{enumerate}[leftmargin=1.6em,itemsep=3pt]
\item \textbf{独立晶格回执绑定了更早的一代。} 独立重放证书 {\footnotesize\texttt{verify/\allowbreak{}out/\allowbreak{}independent\_\allowbreak{}lattice97\_\allowbreak{}upperbound\_\allowbreak{}n100.json}} 记录了
{\footnotesize\texttt{input\_\allowbreak{}sha256 = 7e732e189a0d69e30184dbd053ae2eda6ee588f0418d49553d71004e5f4ff475}} 与
{\footnotesize\texttt{proof\_\allowbreak{}sha256 = a0454013448a2612475975b539d64fa722ab9e69434ca0d1d913bfba03d57d4b}}，而当前的主证书文件哈希为
{\footnotesize\texttt{7a98772a8fb4c608ca32ffe6489bf1f09006dd3c59a4322125beca8c1e72cd7d}}，当前的证明文件哈希为
{\footnotesize\texttt{73d39456b70fdcae41ede5718a608d139661a6c7f2a626cf95b4e348be19600b}}。它的
{\footnotesize\texttt{checker\_\allowbreak{}sha256 = 9655d038a2194ebf2f849ad550ccc645f10c5f7516270d7dfe47267cb23ac878}} \emph{确实}
与当前的独立脚本相符。所以检查器绑定（binding）是当前的，而输入与证明绑定是过期的：独立重放是针对更早的一对证书/证明做的。这并不使对\emph{断言}的独立重放失效（在当前输入上重跑时那 $17$ 项检查通过），但按所存储的样子，该回执不是同代绑定，并且不得被当作同代绑定来呈现。
\item \textbf{一条被引用的证书路径不在镜像中}（见第七部分的产物覆盖子节）。
\item \textbf{一个被记录的常数不在已发布的树中。} 尾部预算使用了一个从被吸收包中导入的量；该预算在
\texttt{verify/\allowbreak{}n100\_\allowbreak{}candidate\_\allowbreak{}budget\_\allowbreak{}222.py} 中的本地重新表述使用
$A=129903810568/10^{11}$，把它连同冻结的斜率（slope）、罚系数（penalty coefficient）与被标记常数一起代入，恰好重现两个存储裕量（见附录中的预算复核）。该常数与 $3\sqrt3/4$ 的一个向外有理包络相同，这一点是从算术得出的\emph{推断}，而不是对该包的引用。
\end{enumerate}

\section{外部材料与吸收边界}

\begin{center}
\footnotesize
% 第 1 列需容纳 56 字符的 \texttt 仓库路径加 40 位十六进制 SHA，0.38\textwidth 不够；
% 下面的宽度为唯一的间隔留出 2\tabcolsep = 12pt。
\begin{tabular}{@{}>{\raggedright\arraybackslash}p{0.45\textwidth}>{\raggedright\arraybackslash}p{\dimexpr0.55\textwidth-2\tabcolsep\relax}@{}}
\toprule
来源 & 角色与边界 \\
\midrule
\texttt{cjt42554-\allowbreak{}crypto/\allowbreak{}disk100-\allowbreak{}covering-\allowbreak{}certificates}
@ \texttt{67ebf1c5776f72a7b7786a8655eed5463429ae55} &
提供被重放的 $100$ 中心构型与六分支骨架；$625$ 个文件，
$7\,269\,413$ 字节；完整性块自身的 scope 是「仅发布完整性；数学结果需要重放」。它的 \texttt{not\_\allowbreak{}claimed} 列表是明确的：全局最优、$\rD(100)$ 的极小多项式、形式化证明助手的几何 \\
\texttt{VonEquinox/\allowbreak{}DiskCoveringSolve} @ \texttt{0895849} &
读取、重放并独立交叉核对：$n=19$ 双链与 $n=20$ 链在隔离副本中重跑，退出码 $0$；$n=11..18$ \emph{未}被重放；本项目不重新分发任何第三方源码 \\
\texttt{VonEquinox/\allowbreak{}DiskCover} & HTTP 请求返回 $404$，Git 报告「Repository not found」；其中没有任何内容被归属或采纳 \\
\bottomrule
\end{tabular}
\end{center}

吸收报告陈述了每个可复用成分的转移条件 --- 完整的拓扑剥离与安全剪枝；精确的力、Farkas 与能量证书；带连续域证明树的候选锚隔离；线性规模的半代数覆盖模型；带实量词消去的单四次平方松弛；以及带隔离的精确端点多项式。在每种情形下，条件都是 $n=100$ 必须重新证明族范围、必要条件与递归完备性：$n=19/20$ 的证书不能当作 $n=100$ 的证书复用。

\section{负对照（negative control）清单}

这里只有在某份证书的破坏性测试被拒绝时，该证书才被接受。正在使用的电池组跨各条路线包括：一个被放大的目标圆盘（仅作前提（precondition），明确不是对所交付断言的未覆盖见证（witness））；一个针对放大圆盘变体的显式未覆盖点；一个缩小与一个放大的覆盖半径（covering radius）；一个错误的度量（六边形内积中缺失的因子 $3$）；一个旋转过的与一个膨胀过的六边形；一次被截断的邻居搜索；一个被错误归属的邻居；两个被移除的相切中心；一个被移除的内部中心；把 $97$ 中心构型当作 $100$ 点断言提交；移除单个壳层中心（覆盖仍然成立 --- 只有基数被拒绝）；一个被破坏的 dart 配对；一个被破坏的面置换；一个被声明为不活跃的支撑；一次断言被禁用的运行；一个不匹配的旧半径；十个被破坏的有理载荷；以及对本地 Krawczyk 见证的 $8$ 项篡改测试（零逆矩阵行、零 Cholesky 变换、重复的主元列、非二进制的逆分母、错误的见证维数、取负的拉格朗日 Hessian、见证源 SHA 不匹配、被改动的几何多项式源码）。此外还有两个入口点拒绝在 \texttt{python -\allowbreak{}O} 下运行。

\section{环境}

已发布的运行使用 Windows 11 与 Anaconda Python 3.13.5（\texttt{sympy} 1.13.3、
\texttt{mpmath} 1.3.0、\texttt{gmpy2} 2.2.1、\texttt{python-\allowbreak{}flint} 0.9.0），文档构建使用 \texttt{tectonic}。判定层只需要标准库：\texttt{fractions}、
\texttt{decimal}、\texttt{math}、\texttt{json}、\texttt{hashlib}、\texttt{pathlib}、
\texttt{itertools}。符号工作使用 \texttt{sympy} 做恒等式检查与因式分解，从不把它当作数值证书的替代品。服务器端的发现层在一个 \texttt{crystalx} 环境中运行，未做任何环境安装；六候选服务器搜索在 $15.746$ s 内完成 $916$ 次目标函数求值，并取回 $25$ 个文件，其 SHA-256 全部验证通过。唯一一次外部项目重放只需要一个 Windows 兼容层（路径分隔符、可执行文件名、LF 输出）；数学条件、常数、区间与计数均未改变。

\part*{第八部分：公开义务，不作软化}
\addcontentsline{toc}{part}{第八部分：公开义务，不作软化}
\stepcounter{partanchor}\setcounter{section}{0}

\section{确切地已确立的内容}

为免歧义，本卷宗所能支持的最强陈述是：

\begin{enumerate}[leftmargin=1.6em,itemsep=3pt]
\item \textbf{一个上界。} 存在一个显式的 $100$ 中心有理解，其覆盖半径（covering radius）至多为 $U_\star=1154469548855829/10^{16}=0.1154469548855829$，由
$138\,437$ 个精确整数对象认证（$128\,436$ 个内部三元组、$9\,900$ 个边界对、$100$
个对径对、$1$ 个例外情形），并带有精确的覆盖平方半径界
$1332799939235381342323027277241/10^{33}=U_\star^2$，把参考值
$115446954885583/10^{15}$ 恰好改进 $1/10^{16}$。
\item \textbf{一个下界。} $B31$--$B35$ 在 $R=2831/25000$ 处由精确证书（certificate）排除；
$B\ge37$ 在 $R=2831/25000$ 处以截距（intercept）$19/2500$ 排除；$B=36$ 在
$R_\star=22647/200000$ 处以截距 $939/125000$ 在 $1\,841\,089$ 个有理盒上排除。半径
单调性随后同样给出 $R_\star$ 处的 $B31$--$B35$ 与 $B\ge37$，故
\[
\rD(100)>22647/200000=0.113235 .
\]
\item \textbf{不依赖构型的界。} $\rD(100)\ge2787/25000=0.11148$（长弦），并有
$\ge1111720153/10^{10}$（加强包络）与 $\ge554570693/5000000000$（整数边计数）
同样成立。
\item \textbf{一个精确的障碍。} 在冻结的 capped-pair 常数下，$B=36$ 松弛在
$R=2831/25000$ 处于有理点 $x=y=9921033/12500000$ 上同时超过旧截距与面积所允许的最大截距。这是关于松弛的陈述，而不是关于覆盖的陈述。
\end{enumerate}

其余一切都或者是条件性的、诊断（diagnostic）性的、接口层面的，或者是公开的。

\section{公开义务}

{\footnotesize
\begin{longtable}{@{}>{\raggedright\arraybackslash}p{0.30\textwidth}>{\raggedright\arraybackslash}p{0.58\textwidth}@{}}
\toprule
义务（obligation） & 缺失内容的精确陈述 \\
\midrule
\endfirsthead
\toprule
义务（obligation） & 缺失内容的精确陈述 \\
\midrule
\endhead
\bottomrule
\endlastfoot
$\rD(100)$ 的精确值 & 未获得。被接受的区间宽度约为 $0.0022$；没有任何论证把它闭合。\texttt{global\_\allowbreak{}optimality\_\allowbreak{}proved = false} \\
$\rD(100)$ 的极小多项式 & 未获得：不是系数，不是次数，不是不可约性，也不是有理的根隔离区间。存在的是\emph{某一个构型}覆盖半径的极小多项式
$67t^2-1$，那是一个不同的对象，并且被明确地如此划定范围 \\
闭合成单个代数数 & 两个端点是不相同的数；没有任何东西把理想包络的零点与最优半径等同起来，而包络的零点明确地是关于该方法极限的陈述，而不是关于 $\rD(100)$ 的陈述 \\
完整的内部映射枚举 & 唯一已知的必要计数定义域是被接受的归一化几何归约下得到的 $87$ 个 $(m,B)$ 对，其中
$m=96..100$（$96:B32..35$；$97:B29..42$；$98:B28..46$；$99:B28..50$；$100:B28..54$）。它不是完整的拓扑枚举，而且存活下来的族中仍含有不可行的图 \\
$m<100$ 与其他图 & 计数桥在其被重放的各半径处排除了 $m\le99$，但对 $m<100$、对其他图族以及对退化的普遍处理尚未闭合 \\
退化 & 重合的中心与冗余的圆盘由归一化处理；其余的退化目录（相切接触、共线中心、边界多重活跃）在各个分支（branch）内部由显式论证处理，而不是由单一的、统一的机器检查处理 \\
$B=36$ 不等式 & 直接的配对/角度松弛无法闭合 $0.11324$；一个共同端点的透镜约束或一个真正耦合的多边平均不等式已被认定为缺失的成分，但它没有通过的全局证书 \\
可续跑的尾部基础设施 & $0.113236$ 与 $0.113237$ 的尝试是停在墙钟时间上，而不是停在数学上；一个带缓存的逆三角与支撑区间、并按对角/非对角分片的可续跑盒队列已被设计出来，但尚未交付 \\
精确值后端 & 被吸收的平方松弛 / 实量词消去理论是唯一已知的通向精确值的可终止路线，但它没有 $n=100$ 的后端；被编译系统的零点集处处奇异，因此常规的 Krawczyk 盒无法替代完整的 QE 后端，而变量预算（$605$ 个几何 $+$ $1318$ 个松弛，$1922$ 个被消去）是一个预算，而不是一个运行时间 \\
遗留源码激活 & 全部 $174$ 个遗留条件源码返回
\texttt{PRECONDITION\_\allowbreak{}MISSING}；依赖门是 \texttt{UNDECIDED}。「每一个历史变体都可直接调用并被逐一重新计时」\emph{未}达成 \\
输出隔离欠账 & 计时、组合/缩放与共享支撑这三个入口点仍然缺少统一的、可重复输出的包装器；整合报告把这列为尚欠的工作 \\
几何形式化 & 不存在 Lean/Coq 内核验证。每一条全称陈述都是解析散文加精确有限算术；解析/几何之间的桥是一个经人工审阅的依赖，在每份证书中都记录为 \texttt{machine\_\allowbreak{}verified\_\allowbreak{}geometry = false} \\
未发布的常数 & 至少有一个尾部常数是从一个不属于镜像的包导入的；它的值在本地被重新表述并数值地重现（附录~\ref{app:budget}），但把该常数与 $3\sqrt3/4$ 的一个向外包络等同起来是一个推断 \\
记号欠账 & 报告没有给出计数参数 $M$ 的自足定义；
它只通过整数超出恒等式与存活剖面出现。读者必须从证书中把它重建出来，而本专著指出这一点，而不是发明一个定义 \\
\end{longtable}}


\section{被撤回或被更正的断言}

一份研究记录只有在保留其撤回时才可信。以下断言曾被提出，随后被撤回或被更正，而这些更正也是记录的一部分：

\begin{itemize}[leftmargin=1.4em,itemsep=3pt]
\item \textbf{$4\arcsin(r/2)$ 弧界。} 错误；$\partial D\cap B(c_i,r)$ 的正确最大角测度是
$2\arcsin r$，在 $\|c_i\|=\sqrt{1-r^2}$ 处取到。错误的常数
在 $r=1$ 附近无法推出 $<\!\pi$，因而不可用。
\item \textbf{边界环带「全局无增益」的接受。} 已作废：一个错误的内面过滤器使待定集为空，而独立脚本复制了同一个错误。退出码 $0$
不是覆盖域正确的证据。
\item \textbf{stage-47 v2 对全部 $4776$ 个构型的排除。} 已否决。
\item \textbf{角能量惯性断言 $+25/-6$。} 已撤回：对真实数据做 CAS 得到 $32\times32$ 角坐标中的一维负方向（外加一个零 = 全局旋转），投影到 $31$ 维间隙坐标后仍是一维负方向。不再作任何惯性断言。
\item \textbf{「径向壳 LP 给出接近 $92.56$ 的已认证覆盖数界。」}
已更正：采样解在 $4001$ 点细网格上不可行（覆盖 $0.99903<1$，
对偶 $1.00317>1$）；唯一被认证的陈述是 $r=L$ 处的粗括界
$[1/r^2,\ \pi/\mathrm{lens}(1,r)]=[80.911,\,165.734]$。更强的说法「这条路线已关闭／证书不可能」也被撤回。
\item \textbf{「构型一侧仍有 $5$--$6\%$ 的余地。」} 已撤回：小 $n$ 的最优解是环型族，而最好的 $n=100$ 构型是晶格型族，因此它们不可比较；晶格族的比值 $c(n)$ 为 $2.13$（$n=40$）、$1.92$（$50$）、$1.35$（$70$）、
$1.14$（$100$）、$1.12$（$140$）、$0.82$（$200$）。
\item \textbf{stage-142 的发现层数字}（耳近似 $1.332634$ 对真值
$1.42016417919847$，以及一处 $K$-目标勘误）不得作为证书值引用。
\item \textbf{stage-44 的第一版}（$4496$ 次收缩 / $597$ 次关联）已被
$2248/562$ 取代；\textbf{stage-45/46} 使用了错误的外 dart 所有者；依赖审计（audit）中一处历史脚本哈希不匹配是被\emph{诊断}出来的，而不是被悄悄修复的。
\end{itemize}

\section{对差距的诚实解读}

两个数字勾勒出局面。不依赖构型的最好下界是 $0.11148$，它比引擎所接受的界 $0.113235$ 低
$1.55\%$；而被接受的区间是
$[0.113235,\,0.1154469548855829]$，相对宽度约为 $1.92\%$。因此引擎比任何单个可用的解析不等式都强得多，而剩下的工作不是更多精度或更多盒子的
问题：它是针对一个分支、在一个半径上的一个不等式，外加 --- 对于精确值而言 --- 一个在 $n=100$ 处尚不存在的枚举与消去后端。

同样重要的是说明当前的界\emph{不}对什么敏感。增加盒子、提高截距、复用 $B\ge37$ 的截距，或者在冻结族内同步重调斜率与罚系数，都会因可证明的理由而失败（第五部分）。该界也不受算力限制：$B=36$ 的候选运行只花 $130$ s 的报告时间，而整个引擎中最大的单项开销是尾部细分。限制它的是：当前局部配对表达式在障碍点（obstruction point）处确实大于 $B=36$ 面积预算所能支付的量。

\appendix
\part*{附录：证据台账、词汇与精确算术}
\addcontentsline{toc}{part}{附录：证据台账、词汇与精确算术}
\stepcounter{partanchor}\setcounter{section}{0}

\section{证据台账（evidence ledger）：断言 $\leftrightarrow$ 产物（artifact）}

{\footnotesize
\begin{longtable}{@{}@{}>{\raggedright\arraybackslash}p{0.28\textwidth}>{\raggedright\arraybackslash}p{0.15\textwidth}>{\raggedright\arraybackslash}p{0.40\textwidth}@{}}
\toprule
断言 & 状态 & 产物 \\
\midrule
\endfirsthead
\toprule
断言 & 状态 & 产物 \\
\midrule
\endhead
\bottomrule
\endlastfoot
半径 $\le0.1154469548855829$ 的 $100$ 圆盘构型 & 已重放 &
{\footnotesize\texttt{verify/\allowbreak{}out/\allowbreak{}external-\allowbreak{}disk100-\allowbreak{}bounds-\allowbreak{}67ebf1c.json}} \\
该构型相对参考界的严格裕量 & 已重放 &
{\footnotesize\texttt{verify/\allowbreak{}out/\allowbreak{}external-\allowbreak{}disk100-\allowbreak{}margin-\allowbreak{}67ebf1c.json}} \\
$115447/10^{6}$ 处的独立四叉树交叉核对 & 已重放 &
\texttt{verify/\allowbreak{}out/\allowbreak{}external-\allowbreak{}disk100-\allowbreak{}local-\allowbreak{}quadtree/\allowbreak{}replay.json}（父仓库） \\
孤立的 KKT 根，模旋转的严格局部极小性 & 已重放 &
{\footnotesize\texttt{verify/\allowbreak{}out/\allowbreak{}external-\allowbreak{}disk100-\allowbreak{}strict-\allowbreak{}local-\allowbreak{}67ebf1c.json}}（父仓库） \\
晶格候选满足 $\alpha^2=1/67$ & 独立重放 &
{\footnotesize\texttt{verify/\allowbreak{}out/\allowbreak{}independent\_\allowbreak{}lattice97\_\allowbreak{}upperbound\_\allowbreak{}n100.json}} \\
$B31$ 排除 & \texttt{PASS} & \texttt{verify/\allowbreak{}out/\allowbreak{}n100-\allowbreak{}B31-\allowbreak{}at-\allowbreak{}11324.json} \\
$B32$ 排除（$M\le25$ 与 $M\ge26$） & \texttt{PASS} &
\texttt{verify/\allowbreak{}out/\allowbreak{}n100-\allowbreak{}B32-\allowbreak{}complete-\allowbreak{}149.json} \\
$B33$ 排除（低 $M$、相干 deep-$5$、$108$ 个环模型） & \texttt{PASS} &
\texttt{verify/\allowbreak{}out/\allowbreak{}n100-\allowbreak{}B33-\allowbreak{}complete-\allowbreak{}167-\allowbreak{}v2.json} \\
$B34$ 排除（$23$ 个无弦 $+$ $47$ 个单弦） & \texttt{PASS} &
\texttt{verify/\allowbreak{}out/\allowbreak{}n100-\allowbreak{}B34-\allowbreak{}complete-\allowbreak{}180.json} \\
$B35$ 排除（低 $M$、$M25$、$M26$、$M\ge27$） & \texttt{PASS} &
\texttt{verify/\allowbreak{}out/\allowbreak{}n100-\allowbreak{}B35-\allowbreak{}complete-\allowbreak{}203.json} \\
$R_\star$ 处的 $B\ge36$ & \texttt{PASS} & \texttt{verify/\allowbreak{}out/\allowbreak{}n100-\allowbreak{}tail-\allowbreak{}candidate-\allowbreak{}224.json} \\
$R=2831/25000$ 处的 $B\ge37$ & \texttt{PASS} &
\texttt{verify/\allowbreak{}out/\allowbreak{}n100-\allowbreak{}tail-\allowbreak{}B37plus-\allowbreak{}210-\allowbreak{}a.json} \\
候选标量预算守恒 & \texttt{PASS} &
\texttt{verify/\allowbreak{}out/\allowbreak{}n100-\allowbreak{}candidate-\allowbreak{}budget-\allowbreak{}222.json} \\
$R_\star$ 处六分支（branch）闭合 & \texttt{PASS} &
\texttt{verify/\allowbreak{}out/\allowbreak{}n100-\allowbreak{}candidate-\allowbreak{}complete-\allowbreak{}225.json} \\
无待定分支的六分支台账 & \texttt{PASS} &
\texttt{verify/\allowbreak{}out/\allowbreak{}n100-\allowbreak{}six-\allowbreak{}branch-\allowbreak{}ledger-\allowbreak{}226.json} \\
$0.11324$ 已证明 & \textbf{未证明} &
{\footnotesize\texttt{verify/\allowbreak{}out/\allowbreak{}n100-\allowbreak{}tail-\allowbreak{}support-\allowbreak{}obstruction-\allowbreak{}206.json}} \\
$0.113236$ / $0.113237$ 已证明 & \textbf{未证明} &
{\footnotesize\texttt{docs/\allowbreak{}result-\allowbreak{}n100-\allowbreak{}stage229-\allowbreak{}current-\allowbreak{}best.md}} \\
整数边计数界 $554570693/5000000000$ & 已接受（解析 $+$ 算术） &
{\footnotesize\texttt{verify/\allowbreak{}out/\allowbreak{}strengthened\_\allowbreak{}lowerbound\_\allowbreak{}n100.json}} \\
加强包络界 $1111720153/10^{10}$ & 已接受 &
\texttt{verify/\allowbreak{}out/\allowbreak{}n100-\allowbreak{}algebraic-\allowbreak{}bounds-\allowbreak{}bundle.json} \\
长弦界 $2787/25000$ & 已接受 &
\texttt{verify/\allowbreak{}out/\allowbreak{}n100-\allowbreak{}longchord-\allowbreak{}bundle-\allowbreak{}frozen.json} \\
面积界 $1/10$ & 已接受 & \texttt{verify/\allowbreak{}out/\allowbreak{}lowerbound\_\allowbreak{}n100.json} \\
$R=153/1250$ 处的连续覆盖 & \texttt{PASS} &
\texttt{verify/\allowbreak{}out/\allowbreak{}coverage\_\allowbreak{}n100.json} \\
独立连续覆盖重放 & 已验证 &
\texttt{verify/\allowbreak{}out/\allowbreak{}independent\_\allowbreak{}coverage\_\allowbreak{}n100.json} \\
覆盖对照电池 & $12/12$ 通过 &
\texttt{verify/\allowbreak{}out/\allowbreak{}coverage\_\allowbreak{}controls\_\allowbreak{}n100.json} \\
角能量路线不可用 & 反例 & {\footnotesize\texttt{verify/\allowbreak{}out/\allowbreak{}n100-\allowbreak{}angular-\allowbreak{}energy-\allowbreak{}obstruction-\allowbreak{}audit.json}} \\
全局 Dirichlet 能量转移失败 & \texttt{H0\_\allowbreak{}PSD = false} &
\texttt{runtime/\allowbreak{}n100-\allowbreak{}vonequinox-\allowbreak{}transfer/\allowbreak{}certificate.json} \\
一致见证对偶无法达到 $>100$ & 已反驳 &
{\footnotesize\texttt{verify/\allowbreak{}out/\allowbreak{}n100-\allowbreak{}witness-\allowbreak{}dual-\allowbreak{}obstruction.json}} \\
边界归属连通性剪枝 & 已反驳（反例） &
{\footnotesize\texttt{verify/\allowbreak{}out/\allowbreak{}n100-\allowbreak{}disconnected-\allowbreak{}boundary-\allowbreak{}owner-\allowbreak{}counterexample.json}} \\
\end{longtable}}


\section{证书词汇}

{\footnotesize
\begin{longtable}{@{}@{}>{\raggedright\arraybackslash}p{0.28\textwidth}>{\raggedright\arraybackslash}p{0.60\textwidth}@{}}
\toprule
字段 & 含义 \\
\midrule
\endfirsthead
\toprule
字段 & 含义 \\
\midrule
\endhead
\bottomrule
\endlastfoot
\texttt{status} & \texttt{PASS}, \texttt{PASS\_\allowbreak{}PREMISES}, \texttt{NOT\_\allowbreak{}CLOSED},
\texttt{PASS\_\allowbreak{}OBSTRUCTION\_\allowbreak{}TO\_\allowbreak{}OLD\_\allowbreak{}RELAXATION}, \texttt{PASS\_\allowbreak{}REDUCTION} \\
\texttt{scope} & 状态在其中成立的确切定义域；一个 \texttt{PASS} 在其之外不作出任何断言 \\
\texttt{radius} & 证书得以建立时所处的候选半径 \\
\texttt{required\_\allowbreak{}center\_\allowbreak{}norm\_\allowbreak{}squared\_\allowbreak{}upper} & $1-R^2$ 的向外有理包络 \\
\texttt{minimum\_\allowbreak{}area\_\allowbreak{}margin} / \texttt{minimum\_\allowbreak{}reported\_\allowbreak{}positive\_\allowbreak{}margin} & 该分支中最小的精确
有理裕量 \\
\texttt{physical\_\allowbreak{}edge\_\allowbreak{}pairs} & 反向半边缘项相消的共享物理边的数目 \\
\texttt{hinge\_\allowbreak{}payment}, \texttt{variance\_\allowbreak{}payment} & 该分支是否花掉其唯一的铰链预算或
方差预算 \\
\texttt{geometric\_\allowbreak{}dependencies} & 该分支所继承的散文层面的解析引理 \\
\texttt{machine\_\allowbreak{}verified\_\allowbreak{}geometry} & 此处每份证书中都是 \texttt{false}：有限算术经机器检查，几何归约则不然 \\
\texttt{global\_\allowbreak{}optimality}, \texttt{global\_\allowbreak{}bound\_\allowbreak{}improved} & 每份分支
证书中都是 \texttt{false}；仅在装配证书中为 \texttt{true} \\
\texttt{bindings} & 源码/输入 SHA-256 映射 \\
\texttt{not\_\allowbreak{}claimed} & 该产物不作断言的显式列表 \\
\end{longtable}}


\section{$B=36$ 障碍与标量预算，精确地}
\label{app:budget}
\label{sec:budget-recheck}

\subsection{障碍}

在 $R=2831/25000$ 处，取 $\lambda=3/50$、$\mu=92531/10^{7}$、被标记常数 $K=973/3125$ 与旧
截距 $L=1501/200000$，点 $x=y=9921033/12500000$ 满足
\[
\frac{3754996055077}{500000000000000}\ \le\ \text{pair}\ \le\
\frac{3754996055079}{500000000000000},
\]
故
\[
\text{pair}-L\ \ge\ \frac{2496055077}{500000000000000}\approx4.99211\times10^{-6},
\]
而 $B=36$ 的面积预算最多允许
\[
L_{\max}=\frac{27019909263463328303182469778967931605940457120574226219}
{3600000000000000000000000000000000000000000000000000000000},
\]
给出超出量
\[
\frac{16062333091071696817530221032068394059542879425773781}
{3600000000000000000000000000000000000000000000000000000000}>0 .
\]
同一点处的直接项分解是
\[
\lambda\text{-distance}\in\Bigl[\frac{1666193587123}{10^{15}},\frac{416548396781}{250000000000000}\Bigr],
\quad
\tfrac12H\in\Bigl[\frac{3338279809719}{500000000000000},\frac{83456995243}{12500000000000}\Bigr],
\]
\[
\mu\text{-angle penalty}\in\Bigl[\frac{416380548203}{500000000000000},\frac{832761096407}{10^{15}}\Bigr],
\qquad
\frac{\text{distance}}{R^2}=\frac{27769893118719}{12823297600000}.
\]

\subsection{所实现的预算，以及一次独立复核}

候选脚本所使用的标量预算是
\[
\text{margin}=(1-\mu)\pi_{\text{lo}}-\bigl((2m-2-B)A+BK\bigr)R^2-BL,
\qquad
\text{gain}=(A-K)R^2-L,
\]
其中 $m=100$、$A=129903810568/10^{11}$。代入冻结常数：
{\footnotesize
\begin{longtable}{@{}llll@{}}
\toprule
候选 & $R$ & $L$ & 裕量 / 增益（十进制） \\
\midrule
\endfirsthead
\toprule
候选 & $R$ & $L$ & 裕量 / 增益（十进制） \\
\midrule
\endhead
\bottomrule
\endlastfoot
stage 221/222 & $113237/10^{6}$ & $75097/10^{7}$ &
$4.906445222464608\times10^{-7}$ \;/\; $5.1549192230158705\times10^{-3}$ \\
candidate 224 & $22647/200000$ & $939/125000$ &
$1.8087101250224213\times10^{-5}$ \;/\; $5.152171860143971\times10^{-3}$ \\
\end{longtable}}

两者都精确重现了所存储的证书值，公式正是以此被检查的。所存储的精确裕量是
\[
\text{candidate 222: }\quad \frac{49064452224646078469778967931605940457120574226219}{10^{56}},
\]
\[
\text{candidate 224: }\quad \frac{1808710125022421182469778967931605940457120574226219}{10^{56}},
\]
而所存储的每胞增益是 $64436490287698376249/12500000000000000000000$ 与
$2576085930071985489/500000000000000000000$。

\section{合并的分支裕量}

裕量按半径分别列出；它们在各自半径上计算，不得跨半径比较或相加。

\begin{itemize}[leftmargin=1.4em,itemsep=2pt]
\item $B31$：在 $R=1132301/10^{7}$ 处为
$\frac{3791306959813284956819}{10^{24}}$；在 $R=2831/25000$ 处为
$\frac{3205113070935301276991}{10^{24}}$；在 $R_\star$ 处为继承。
\item $B32$：在 $R=1132301/10^{7}$ 处为
$\frac{40052428178643350631}{125000000000000000000000}$；在 $R=2831/25000$ 处为正，其中 $73$
个原始行至少为 $0.00007800$，$8$ 个新行大于 $0.00012768$；在 $R_\star$ 处为
继承。
\item $B33$：在 $R=1132301/10^{7}$ 处为
$\frac{141216899554184842489}{10^{24}}$；在 $R=2831/25000$ 处至少为
$\frac{1471448836581788939}{500000000000000000000000}$；在 $R_\star$ 处为继承。
\item $B34$：在 $R=1132301/10^{7}$ 处为
$\frac{57252253086154121097}{250000000000000000000000}$；在 $R=2831/25000$ 处至少为
$\frac{190482881108285619519923}{292968750000000000000000000}$；在 $R_\star$ 处为继承。
\item $B35$：在 $R=1132301/10^{7}$ 处为
$\frac{37566017111533099749}{500000000000000000000000}$；在 $R=2831/25000$ 处至少为
$\frac{169690032082888672651865333931}{625000000000000000000000000000000}$；在 $R_\star$ 处为继承。
\item $B\ge36$：在 $R=1132301/10^{7}$ 处为
$\frac{33605092614133890926309778967931605940457120574226219}{10^{56}}$；在 $R=2831/25000$ 处，
$B=36$ 受阻，而 $B\ge37$ 给出
$\frac{166438291677217339982469778967931605940457120574226219}{10^{56}}$；在 $R_\star$ 处，$B=36$
给出
$\frac{1808710125022421182469778967931605940457120574226219}{10^{56}}$。
\end{itemize}

$B31$ 在 $R=11323/100000$ 处的裕量为 $\frac{3797228001573004612011}{10^{24}}$。


\section{术语表（glossary）}

{\footnotesize
\begin{longtable}{@{}>{\raggedright\arraybackslash}p{0.26\textwidth}>{\raggedright\arraybackslash}p{0.62\textwidth}@{}}
\toprule
术语 & 含义 \\
\midrule
\endfirsthead
\toprule
术语 & 含义 \\
\midrule
\endhead
\bottomrule
\endlastfoot
边界弧（boundary arc） & $\partial D$ 的极大正测度连通段，其归属者几乎处处（a.e.\ ）为常数；
其计数为 $B$ \\
capped & 传播达到其轮次上限而未达到不动点；被保留，未被判定 \\
capped-pair 松弛 & 用于尾部分支 $B\ge36$ 的局部对损失 \\
证书（certificate） & 一份带有状态、scope、计数与哈希的结构化机器产物 \\
核心（core） & 剥离外边界带之后剩下的诱导圆盘 \\
深层 / 环顶点（deep / ring vertex） & 严格位于核心内部的顶点，或分别地位于其边界上的顶点 \\
耳（ear） & 可在保持星形性质的同时被删除的核心面块 \\
坏面（bad face）($q1$) & 被计为亏缺的那一类被标记面 \\
铰链 / 方差预算（hinge / variance budget） & 分支可以花费一次的那个单一标量预算 \\
单调性继承（monotonicity inheritance） & 把在较大半径处证明的排除用于较小半径 \\
障碍点（obstruction point） & 松弛超过其允许预算的精确有理点 \\
向外舍入（outward rounding） & 区间端点朝远离真值的方向舍入 \\
见证（witness） & 用于导出一个必要条件的一个点或一个点集 \\
\end{longtable}}


\section{本文档的审计（audit）记录}

在准备本专著期间执行了以下检查，其结果在正文中报告，而不仅在本附录中：

\begin{enumerate}[leftmargin=1.6em,itemsep=3pt]
\item \texttt{python reproduce.py} --- 退出码 $0$；被接受的界 $0.113235$；$0.11324$ 作为
未闭合被保留；打印四个证书哈希（第七部分，入口点子节）。
\item \textbf{产物覆盖审计} --- 镜像（mirror）自身文档引用的 $100$ 条路径中，
$19$ 条存在于镜像中，$81$ 条不存在；这 $81$ 条中有 $80$ 条存在于父仓库；
\texttt{n100-\allowbreak{}public/\allowbreak{}proofs/\allowbreak{}} 缺失，而被引用的证明存在于父仓库，
其 SHA-256 与证书的 \texttt{proof\_\allowbreak{}sha256} 相符（第七部分，产物覆盖子节）。
\item \textbf{绑定（binding）审计} --- 独立晶格回执绑定了一个过期的输入与证明
（\texttt{7e732e18\ldots} 与 \texttt{a0454013\ldots}），而其检查器绑定
（\texttt{9655d038\ldots}）与当前脚本相符（第七部分，绑定审计子节）。
\item \textbf{数值重新推导} --- 标量预算公式在两个候选半径上以精确有理
算术重新求值，并精确重现所存储的裕量与增益
（附录~\ref{app:budget}）。
\item \textbf{分数位数审计} --- $B=36$ 的裕量是 $52$ 位分子除以 $10^{56}$；
障碍的面积允许截距是 $56$ 位分子除以 $36\cdot10^{55}$；
覆盖平方半径界是 $31$ 位分子除以 $10^{33}$，等于 $U_\star^2$。
\item \textbf{移植包检查} --- 分支证书所引用的四个源码存在于父仓库的移植包中，
且不在镜像中。
\end{enumerate}

\begin{remark}
本附录有意把失败与被撤回的条目与已接受的条目并列。一份只会自我确认的记录从未受过检验；上面这份清单的价值在于它显示了该计划自身的机器在何处出错，以及是什么取代了它。
\end{remark}


\end{document}
